\chapter{Sharpness and Optimal Blocking}\label{ch4-sharpness-blocking}\label{ch:sharpness}

\begin{quote}\itshape
This chapter asks how much of the certificate of Chapter~\ref{ch:certification} is unimprovable and how its one free design parameter should be set. Three things are settled. The constant $1$ on the Kolmogorov drift term is sharp: for every calibration law with a continuous distribution function and every threshold procedure there is a deployment law at Kolmogorov distance $\varepsilon$ whose coverage is lower by $\varepsilon$ up to two corrections, one of which is exponentially small for split conformal. The block length that minimises the expected calibration-conditional shortfall over the polynomial mixing class $\Bclass{r}$ is computed exactly, with its constant, and is not the point at which the two terms of the bound balance. And the blocked rate so obtained is not minimax: at fixed confidence the unblocked sample quantile of the whole record attains $N^{-1/2}$ for every $r>1$, which is optimal, so blocking buys the exact Beta law and a failure probability depending only on the assumed class, not rate. The chapter then states, and does not prove, the result that would carry the argument furthest: a mixing-driven lower bound $(N^r\delta)^{-1/(r+1)}$ on the two-sided calibration-conditional coverage error at confidence $\delta$, holding over all threshold procedures on $\Bclass{r}$, together with the regime-renewal construction that should produce it and an exact list of what remains open. Two smaller parts follow: the partially observed drift problem, stated as a target, and the dependent empirical Wasserstein rate, which is specialised from the literature and claimed as nothing more.
\end{quote}

\section{What Chapter~\ref{ch:certification} leaves open}\label{sec:sharp-intro}

Chapter~\ref{ch:certification} produced a certificate. On a record of $N$ periods of hedging errors from a fixed valuation-and-hedging policy, blocked into $n=\lfloor N/(2b)\rfloor$ blocks of length $b$ separated by gaps of length $b$, the split-conformal threshold $\hat q=S_{(k)}$ covers the next gap-separated block with probability at least $1-\alpha$ less a Kolmogorov drift term $\dK(P,Q)$ and a coupling term $(n+1)\beta(b)$, marginally (Theorem~\ref{thm:cov}(i)); and, on an event of probability at least $1-\delta-\delta'-n\beta(b)$, the calibration-conditional coverage is at least $b_{n,k}(\delta)-\dK(P,Q)-\beta(b)/\delta'$ (Theorem~\ref{thm:cov}(ii)). Proposition~\ref{prop:twosided} brackets the same coverage from above under a continuity hypothesis on the calibration score law, and so defines a two-sided calibration-conditional coverage error.

Four questions are left open, and this chapter is their answer.

\paragraph{Is the drift term sharp?} The certificate charges $\dK(P,Q)$ with coefficient $1$. Nothing in Chapter~\ref{ch:certification} shows that a cleverer argument could not charge $\tfrac12\dK$, or charge a smaller distance. In a value-at-risk setting the nearest published bound carries a coefficient $2$ on a total-variation drift term and its author shows the factor $2$ is not removable for that argument \citep[p.~3]{schmitt2026}; the coupling route of Chapter~\ref{ch:certification} charges $\dK$ once, and the worked example there shows the charge is loose by a factor of nearly three at plausible parameters. Whether the looseness is in the bound or in the example is the question. Section~\ref{sec:sharp-thm8} settles it: the bound is exact and the example is favourable. The construction moves the bottom-$\varepsilon$ mass of the calibration law to a far point, and the resulting coverage loss is $\varepsilon$ less two corrections which vanish as the far point recedes and as the calibration sample grows.

\paragraph{How should the block length be chosen?} The block length is not a nuisance parameter. It sets the number of blocks, $n\approx N/(2b)$, and hence the width of the Beta law of calibration-conditional coverage; and it sets the coupling term $(n+1)\beta(b)$, which falls as $b$ grows. Algorithm~\ref{alg:cert} takes $b$ as an input and says nothing about where it should come from. Section~\ref{sec:sharp-block} fixes a mixing class, takes the expected shortfall implied by Theorem~\ref{thm:cov}(ii), and minimises it exactly. The answer is not the balance point of the two terms; the balance point costs a third more than the minimum at the parameters of the worked example.

\paragraph{Is blocking the right design at all?} Half the record is spent on gaps. The blocked route was adopted in Chapter~\ref{ch:certification} because it delivers an exact $\Beta(k,n+1-k)$ law for the calibration-conditional coverage and a drift term charged on a fixed half-line, and because its failure probability depends on the assumed mixing class through one term only. It was not adopted because it is efficient, and Section~\ref{sec:sharp-notminimax} shows that it is not: the ordinary sample quantile of the whole unblocked record has two-sided calibration-conditional coverage error of order $(N\delta)^{-1/2}$ for every $r>1$, and no procedure can do better than $N^{-1/2}$ at fixed confidence, because independent scores lie in every class. The blocked rate is slower by the factor $\sqrt{b^*}$. This is the chapter's least comfortable result and it is stated first among its conclusions.

\paragraph{What is minimax?} The two previous questions are about one procedure and one design. The question a doctorate must answer is about all procedures: over the class $\Bclass{r}=\{\beta(k)\le Ck^{-r}\}$, how small can the two-sided calibration-conditional coverage error be made at confidence $\delta$, by any threshold computed from the record? Section~\ref{sec:sharp-notminimax} answers it at fixed $\delta$, where the answer is $N^{-1/2}$ and the mixing exponent does not appear. Section~\ref{sec:sharp-lower} states the answer at high confidence, where it does appear and where the certificates a risk committee actually signs live: the claim is $(N^r\delta)^{-1/(r+1)}$, produced by a regime-renewal process whose long runs are rare and damaging. That statement is not proved here. It is stated with its construction, its route, and a box listing exactly what is missing, and its plausibility is checked against a published rate: its expected shortfall integrates to $N^{-r/(r+1)}$, which is what \citet[Cor.~2]{barber2026} prove for split conformal marginally.

\paragraph{Map.} Section~\ref{sec:sharp-thm8} proves the sharpness theorem. Section~\ref{sec:sharp-functional} defines the coverage-error functional that the rest of the chapter optimises and explains why it must be two-sided. Sections~\ref{sec:sharp-block} to~\ref{sec:sharp-w1} take the five parts of the blocking theorem in turn: the optimal block length with a full proof; the failure of minimaxity, with a full proof of both the upper and the lower half; the mixing-driven lower bound, with statement, construction, route and gap; the partially observed drift, with statement, route and gap; and the dependent empirical Wasserstein rate, specialised from the literature. Section~\ref{sec:sharp-positioning} positions all of it against the five nearest papers. Section~\ref{sec:sharp-example} gives the numerical illustration, including the sweep over block lengths on the synthetic process of Chapter~\ref{ch:certification}. Section~\ref{sec:sharp-discussion} discusses, and states plainly which parts carry doctoral weight.

The five parts of the blocking theorem are stated separately below, as Theorems~\ref{thm:block-b}, \ref{thm:block-rate}, \ref{thm:block-lower}, \ref{thm:block-drift} and Proposition~\ref{prop:block-w1}, because their status differs: two are proved, one is a target with a route, one is a target with a route and a weaker proved companion, and one is a specialisation of a published result. Read together they are one statement about the calibration-conditional coverage error over a mixing class.

\section{Sharpness of the drift penalty}\label{sec:sharp-thm8}

\subsection{The setting}

Throughout this section $P$ is a law on $\R$ with continuous distribution function $F_P$, and $S_0$ is a deployment score whose law is written $Q$. The object of interest is the coverage that a threshold delivers.

\begin{definition}[Threshold procedure and its coverage]\label{def:sharp-proc}
A \emph{threshold procedure} is a random variable $T$ with values in $\R\cup\{+\infty\}$, defined on a probability space and independent of the deployment score. For a law $R$ on $\R$ write
\[
C_T(R):=\Prob_{S\sim R}(S\le T)=\E\big[F_R(T)\big],
\]
the second equality by independence of $T$ and $S$, with $F_R(+\infty):=1$.
\end{definition}

The split-conformal threshold of Chapter~\ref{ch:certification} is a threshold procedure in this sense once the coupling of Corollary~\ref{cor:coupled} has made the deployment score independent of the calibration record: $T=\hat q^*$ and $C_T(P)=\E F_P(\hat q^*)$, which Corollary~\ref{cor:exact} identifies with $k/(n+1)$ when $F_P$ is continuous. Definition~\ref{def:sharp-proc} deliberately says nothing else about $T$. The theorem below is a statement about every threshold, including thresholds that are not order statistics, thresholds that use side information, and the degenerate threshold $T=+\infty$.

\begin{definition}[Bottom-mass transport]\label{def:sharp-transport}
Let $\varepsilon\in(0,1)$ and let $q_\varepsilon:=\inf\{x:F_P(x)\ge\varepsilon\}$, so that $F_P(q_\varepsilon)=\varepsilon$ by continuity and $F_P(x)<\varepsilon$ for every $x<q_\varepsilon$. For $M>q_\varepsilon$ define the law
\[
Q_M:=P\big|_{(q_\varepsilon,\infty)}+\varepsilon\,\delta_M ,
\]
that is, $P$ with the whole of its mass on $(-\infty,q_\varepsilon]$ removed and replaced by an atom of size $\varepsilon$ at the point $M$.
\end{definition}

\begin{lemma}[Distances to the transported law]\label{lem:sharp-dist}
For every $\varepsilon\in(0,1)$ and every $M>q_\varepsilon$,
\[
\dTV(P,Q_M)=\dK(P,Q_M)=\varepsilon,
\]
and the distribution functions satisfy
\begin{equation}\label{eq:sharp-cdf}
F_P(x)-F_{Q_M}(x)=
\begin{cases}
F_P(x)\in[0,\varepsilon), & x<q_\varepsilon,\\
\varepsilon, & q_\varepsilon\le x<M,\\
0, & x\ge M.
\end{cases}
\end{equation}
\end{lemma}

\begin{proof}
For $x<q_\varepsilon$ the law $Q_M$ places no mass in $(-\infty,x]$, because its mass is confined to $(q_\varepsilon,\infty)\cup\{M\}$; hence $F_{Q_M}(x)=0$ and the difference is $F_P(x)$, which lies in $[0,\varepsilon)$ by the definition of $q_\varepsilon$. For $q_\varepsilon\le x<M$ one has $F_{Q_M}(x)=P((q_\varepsilon,x])=F_P(x)-F_P(q_\varepsilon)=F_P(x)-\varepsilon$. For $x\ge M$ the atom is included and $F_{Q_M}(x)=F_P(x)-\varepsilon+\varepsilon$. This is \eqref{eq:sharp-cdf}, and its supremum is $\varepsilon$, attained on $[q_\varepsilon,M)$, so $\dK(P,Q_M)=\varepsilon$. For total variation, $P$ and $Q_M$ agree on $(q_\varepsilon,\infty)\setminus\{M\}$; the mass transported is $P((-\infty,q_\varepsilon])=\varepsilon$, and $P(\{M\})=0$ because $F_P$ is continuous, so $\dTV(P,Q_M)=\varepsilon$ by Definition~\ref{def:prelim-dK-dTV} of Chapter~\ref{ch:prelim}.
\end{proof}

\subsection{The theorem}

\begin{theorem}[Sharpness of the drift penalty]\label{thm:sharp}
Let $T$ be a threshold procedure in the sense of Definition~\ref{def:sharp-proc} and let $P$ have a continuous distribution function. Fix $\varepsilon\in(0,1)$.
\begin{enumerate}[nosep,label=(\alph*)]
\item (Universal.) For every law $Q$ on $\R$,
\[
C_T(Q)\;\ge\;C_T(P)-\dK(P,Q)\;\ge\;C_T(P)-\dTV(P,Q),
\]
so that $\inf\{C_T(Q):\dK(P,Q)\le\varepsilon\}\ge C_T(P)-\varepsilon$ and the same with $\dTV$ in place of $\dK$.
\item (Attained.) With $Q_M$ as in Definition~\ref{def:sharp-transport},
\begin{equation}\label{eq:sharp-identity}
C_T(Q_M)=C_T(P)-\varepsilon+\varepsilon\,\Prob(T\ge M)+\E\Big[\big(\varepsilon-F_P(T)\big)\ind{F_P(T)<\varepsilon}\Big],
\end{equation}
and consequently
\begin{equation}\label{eq:sharp-bound}
C_T(P)-\varepsilon+\varepsilon\,\Prob(T\ge M)\;\le\;C_T(Q_M)\;\le\;C_T(P)-\varepsilon+\varepsilon\,\Prob(T\ge M)+\varepsilon\,\Prob\big(F_P(T)<\varepsilon\big).
\end{equation}
\item (Split conformal.) If $T=S_{(k)}$ is the $k$-th order statistic of $n$ independent scores with law $P$ and $k=\lceil(1-\alpha)(n+1)\rceil\le n$, then for every $\varepsilon\le 1-\alpha$,
\[
\Prob\big(F_P(T)<\varepsilon\big)=\Prob\big(\Beta(k,n+1-k)<\varepsilon\big)=\Prob\big(\mathrm{Bin}(n,\varepsilon)\ge k\big)\le\exp\big(-2n(k/n-\varepsilon)^2\big),
\]
and $\Prob(T\ge M)\to0$ as $M\to\infty$. Hence for every $\eta>0$ there is $M$ with
\[
C_T(Q_M)\;\le\;C_T(P)-\varepsilon\Big(1-\eta-e^{-2n(k/n-\varepsilon)^2}\Big),
\]
so the coefficient $1$ on $\dK(P,Q)$ in Theorem~\ref{thm:cov} cannot be reduced.
\end{enumerate}
\end{theorem}

\begin{proof}
\emph{(a).} Condition on $T$, which is independent of the deployment score. For a fixed value $t\in\R$ the event $\{S\le t\}$ is a half-line, so
\[
Q\big((-\infty,t]\big)=F_Q(t)\ge F_P(t)-\sup_x\big|F_P(x)-F_Q(x)\big|=F_P(t)-\dK(P,Q),
\]
which is Lemma~\ref{lem:halfline} of Chapter~\ref{ch:certification}; for $t=+\infty$ both sides are $1$ and the inequality is trivial. Taking the expectation over $T$ gives $C_T(Q)\ge C_T(P)-\dK(P,Q)$. The second inequality is $\dK\le\dTV$ (Proposition~\ref{prop:prelim-dK-le-dTV} of Chapter~\ref{ch:prelim}), which holds because half-lines are measurable sets. This is the whole content of (a): the drift term of Theorem~\ref{thm:cov} is nothing but the definition of $\dK$ evaluated at a point that the coupling has frozen.

\emph{(b).} By \eqref{eq:sharp-cdf} and the definition of $q_\varepsilon$, and using $\{T<q_\varepsilon\}=\{F_P(T)<\varepsilon\}$, which is an identity of events because $F_P(x)<\varepsilon$ exactly for $x<q_\varepsilon$,
\[
F_{Q_M}(T)=\big(F_P(T)-\varepsilon\big)+\varepsilon\ind{T\ge M}+\big(\varepsilon-F_P(T)\big)\ind{F_P(T)<\varepsilon}.
\]
The three cases of \eqref{eq:sharp-cdf} are checked directly: on $\{T\ge M\}$ the right side is $F_P(T)$; on $\{q_\varepsilon\le T<M\}$ it is $F_P(T)-\varepsilon$; on $\{T<q_\varepsilon\}$ it is $F_P(T)-\varepsilon+\varepsilon-F_P(T)=0$. The value $T=+\infty$ falls in the first case. Taking expectations gives \eqref{eq:sharp-identity}. For \eqref{eq:sharp-bound}, the last term of \eqref{eq:sharp-identity} is non-negative and is at most $\varepsilon\Prob(F_P(T)<\varepsilon)$, because $0\le\varepsilon-F_P(T)\le\varepsilon$ on the event in question.

\emph{(c).} The identity $F_P(S_{(k)})=U_{(k)}$ in law, for $U_i:=F_P(S_i)$ independent and uniform by the probability integral transform (Lemma~\ref{lem:prelim-pit} of Chapter~\ref{ch:prelim}, $F_P$ continuous), gives $F_P(T)\sim\Beta(k,n+1-k)$ by Lemma~\ref{lem:prelim-order-stat} of Chapter~\ref{ch:prelim}; the same lemma supplies the binomial representation $\Prob(U_{(k)}\le u)=\Prob(\mathrm{Bin}(n,u)\ge k)$, and Lemma~\ref{lem:prelim-hoeffding} of Chapter~\ref{ch:prelim} supplies the Chernoff--Hoeffding bound, applicable because $k/n\ge(1-\alpha)(n+1)/n>1-\alpha\ge\varepsilon$. The order statistic of finitely many real scores is finite almost surely, so $\Prob(T\ge M)\downarrow0$ as $M\uparrow\infty$ by monotone convergence; choose $M$ with $\Prob(T\ge M)\le\eta$ and substitute into \eqref{eq:sharp-bound}. Finally, $\dK(P,Q_M)=\varepsilon$ by Lemma~\ref{lem:sharp-dist}, so the pair $(P,Q_M)$ is admissible in Theorem~\ref{thm:cov} with drift term exactly $\varepsilon$, and the realised coverage loss is $\varepsilon(1-\eta-e^{-2n(k/n-\varepsilon)^2})$. Letting $n\to\infty$ and $\eta\to0$, no coefficient smaller than $1$ can multiply $\dK$ in a bound valid for all $(P,Q)$ and all $n$.
\end{proof}

\subsection{Reading the two correction terms}

The identity \eqref{eq:sharp-identity} is exact, and the two corrections have distinct meanings.

The term $\varepsilon\Prob(T\ge M)$ is the adversary's cost of hiding the mass at a finite point. If the threshold happens to exceed $M$, the transported mass is covered after all and none of the loss materialises. It is a defect of the construction and not of the bound: it can be made arbitrarily small by taking $M$ large, and it vanishes for every threshold that is almost surely finite. It cannot be removed altogether, because the point $+\infty$ is not a point of $\R$ and $\delta_{+\infty}$ is not a probability measure on $\R$. Any construction that tries to place mass beyond every threshold at once leaves the space of laws on $\R$, and with it the hypothesis of Theorem~\ref{thm:cov}.

The term $\varepsilon\Prob(F_P(T)<\varepsilon)$ is the adversary's cost of the threshold falling below the transported mass in the first place. If $T<q_\varepsilon$, the mass on $(-\infty,q_\varepsilon]$ was not being covered before the transport either, so moving it away costs nothing. For split conformal this term is not merely small; it is exponentially small in the number of blocks, because the threshold is a high order statistic and $F_P(\hat q)$ concentrates near $k/(n+1)$.

\begin{example}[Size of the corrections at the parameters of Chapter~\ref{ch:certification}]\label{ex:sharp-corr}
Take the synthetic example of Chapter~\ref{ch:certification}: $n=25$, $k=24$, $\alpha=0.1$, and let $\varepsilon=0.05$, a drift of five percentage points in Kolmogorov distance. Then
\[
\Prob\big(F_P(\hat q)<0.05\big)=\Prob\big(\mathrm{Bin}(25,0.05)\ge24\big)=25(0.05)^{24}(0.95)+(0.05)^{25}=1.4\times10^{-30},
\]
against the Hoeffding bound $\exp(-50(0.96-0.05)^2)=1.0\times10^{-18}$. Both are beyond any relevance. With $M$ chosen so that $\Prob(\hat q\ge M)\le10^{-6}$, the realised coverage loss under $Q_M$ is $\varepsilon$ to six decimal places: the certificate's charge of $\dK(P,Q)$ is exactly right for this $Q$.
\end{example}

\datanote{The two numbers in Example~\ref{ex:sharp-corr} are the closed forms printed there, evaluated in double precision: $25\times0.05^{24}\times0.95=1.416\times10^{-30}$ and $0.05^{25}=2.98\times10^{-33}$, summing to $1.42\times10^{-30}$; and $\exp(-50\times0.91^2)=\exp(-41.405)=1.04\times10^{-18}$. No data enter.}

\begin{remark}[Why the far-mass mixture is weaker]\label{rem:sharp-mixture}
The construction that first suggests itself is the mixture $Q'_M:=(1-\varepsilon)P+\varepsilon\delta_M$, which also moves mass $\varepsilon$ to a far point but takes it proportionally from the whole of $P$ rather than from its bottom. Its distribution function satisfies $F_P(x)-F_{Q'_M}(x)=\varepsilon F_P(x)$ for $x<M$ and $0$ for $x\ge M$, so $\dTV(P,Q'_M)=\varepsilon$ and $\dK(P,Q'_M)=\varepsilon F_P(M^-)\uparrow\varepsilon$ as $M\to\infty$: the distances are the same as in Lemma~\ref{lem:sharp-dist} in the limit. But the coverage is
\[
C_T(Q'_M)=(1-\varepsilon)C_T(P)+\varepsilon\Prob(T\ge M)\;\longrightarrow\;C_T(P)-\varepsilon\,C_T(P),
\]
so the loss is $\varepsilon\,C_T(P)$, not $\varepsilon$. For split conformal $C_T(P)=k/(n+1)\approx1-\alpha$, and the mixture therefore establishes only that the coefficient on $\dK$ is at least $1-\alpha$: at $\alpha=0.1$ it leaves a tenth of the constant unaccounted for, and it is sharp only for degenerate procedures with $C_T(P)=1$, that is, for thresholds that cover everything. The difference between the two constructions is where the mass comes from. The mixture takes a fraction $\varepsilon$ of the mass that was already covered, so a fraction $1-C_T(P)$ of what it moves was worthless to the adversary; the bottom-mass transport of Definition~\ref{def:sharp-transport} takes only mass that was covered with probability $1-\Prob(F_P(T)<\varepsilon)$, which for a high order statistic is essentially one. Choosing the transport to hit exactly the mass the threshold covers is what turns $1-\alpha$ into $1$.
\end{remark}

\begin{remark}[Adversarial, not Le Cam]\label{rem:sharp-lecam}
Theorem~\ref{thm:sharp} is not a minimax statement and its proof is not an application of Le Cam's two-point method \citep[\S2.3, pp.~81--83]{tsybakov2009}. The two-point method is needed when the adversary must choose a law that is simultaneously damaging and hard to distinguish from the truth on the available data, so that the statistician cannot adapt. Here the deployment law $Q$ is never observed: the certificate is issued before the deployment block is realised, and no amount of data about $P$ constrains $Q$ beyond the assumed bound on $\dK(P,Q)$. The relevant quantity is therefore an infimum over a ball, $\inf\{C_T(Q):\dK(P,Q)\le\varepsilon\}$, and Theorem~\ref{thm:sharp} evaluates it by exhibiting the minimiser. No testing problem arises and no total-variation distance between sample laws is needed. Le Cam's method enters this chapter twice, and both times because something \emph{is} observed: in Section~\ref{sec:sharp-notminimax}, where the calibration record is observed and the lower bound must hold over all procedures computed from it, and in Section~\ref{sec:sharp-drift}, where a sample of deployment scores is observed and the size of the drift must itself be estimated.
\end{remark}

\begin{remark}[What sharpness does and does not say]\label{rem:sharp-scope}
Theorem~\ref{thm:sharp} fixes $P$, fixes the procedure, and finds the worst $Q$ in a Kolmogorov ball. It does not say that the certificate's charge is right for the $Q$ a converted book actually faces. The worked example of Chapter~\ref{ch:certification} makes the point in the other direction: there $P$ and $Q$ are normal with equal variance and means differing by $b\mu$, the Kolmogorov distance $0.0391$ is attained at the midpoint of the two means, and the true coverage loss is $0.0142$, because the threshold sits far into the upper tail where the two distribution functions have almost converged. Both statements are true and neither weakens the other. A bound that holds for every $Q$ at a given $\dK$ must be tight at the worst $Q$ and is therefore loose at benign ones. The practical consequence is stated in Section~\ref{sec:sharp-discussion}: if a bank can say more about the shape of its drift than a Kolmogorov bound, the certificate can be sharpened, and Theorem~\ref{thm:sharp} says exactly where the room is, namely in the class of admissible $Q$ and nowhere in the argument.
\end{remark}

\section{The conditional coverage error and why it must be two-sided}\label{sec:sharp-functional}

Sections~\ref{sec:sharp-block} to~\ref{sec:sharp-lower} optimise a functional, and it must be defined before it can be optimised. Two decisions are embedded in the definition: the error is conditional on the calibration record, and it is two-sided.

\begin{definition}[Calibration-conditional coverage error]\label{def:sharp-error}
Let $(S_t)_{t\ge1}$ be a strictly stationary process on the calibration record with continuous marginal distribution function $F_P$, and let $S_{1:N}$ be the record. A \emph{threshold procedure} is a measurable map $T=T(S_{1:N})$ with values in $\R\cup\{+\infty\}$. Its \emph{conditional coverage error} at level $\alpha$ is
\[
\mathcal E(T):=\big|F_P(T)-(1-\alpha)\big| ,
\]
the error in the coverage that a deployment score independent of the record receives from the threshold. Its \emph{coverage error at confidence $\delta$} is
\[
\mathcal E_\delta(T):=\inf\big\{\varepsilon\ge0:\Prob\big(\mathcal E(T)>\varepsilon\big)\le\delta\big\},
\]
and, for the class $\Bclass{r}=\Bclass{r}(C):=\{\text{processes with }\beta(k)\le Ck^{-r}\text{ for all }k\ge1\}$ with $r>1$, the \emph{minimax coverage error} is
\[
\mathcal E^*_\delta(N,r):=\inf_T\ \sup_{\Prob\in\Bclass{r}}\ \mathcal E_\delta(T).
\]
\end{definition}

Three features of the definition need comment.

\paragraph{Why conditional.} The quantity $F_P(T)$ is the coverage delivered by the particular threshold in hand, not averaged over records. It is the quantity Theorem~\ref{thm:cov}(ii) and Proposition~\ref{prop:twosided}(b) bound, and it is the quantity a risk manager needs: having computed this threshold from this record, how much of the deployment law does it cover? The marginal coverage $\E F_P(T)$ is a different object, it is settled for split conformal under $\beta$-mixing by \citet[Cor.~2]{barber2026} without blocking, and its minimax analysis over all procedures runs into the degeneracy described next. Definition~\ref{def:sharp-error} therefore takes the conditional object as the thing to be optimised, and Section~\ref{sec:sharp-positioning} records that the marginal problem is left to the literature that has solved it.

\paragraph{Why two-sided.} If $\mathcal E$ were replaced by the one-sided deficit $(1-\alpha-F_P(T))^+$, the infimum over procedures would be zero for every $N$, $r$ and $\delta$: the threshold $T\equiv+\infty$ has deficit zero, covers everything, and is useless. A minimax statement over all threshold procedures must therefore either charge for over-coverage or impose an efficiency constraint, and charging for over-coverage is the simpler of the two and the one already in the literature: \citet[Thm.~3]{barber2023} bound over-coverage by the same quantity that bounds under-coverage in their Theorem~2, and \citet[Thm.~4]{halkiewicz2026} take the two-sided coverage error itself as the object of a minimax statement. Over-coverage is not a technical device. In the application of this thesis the threshold is a reserve in currency units per block, and a threshold that over-covers is a reserve that is too large: capital held against a loss that will not occur. Remark~\ref{rem:twosided} of Chapter~\ref{ch:certification} makes the same point, and Proposition~\ref{prop:twosided}(c) of that chapter is the upper bound on $\mathcal E$ that the blocked route delivers.

\paragraph{Why the drift and the deployment gap are excluded.} Definition~\ref{def:sharp-error} contains no drift term and no coupling cost for the deployment block. Both are charged in Chapter~\ref{ch:certification} and both are additive there: Theorem~\ref{thm:cov}(ii) subtracts $\dK(P,Q)$ and $\beta(b)/\delta'$ from a bound on $F_P(\hat q)$, and Proposition~\ref{prop:twosided}(b) adds them. Excluding them here is legitimate precisely because Theorem~\ref{thm:sharp} has just shown the drift charge to be exact: nothing about the design of the threshold can reduce it, so it does not interact with the optimisation. What remains is the pure calibration-side error, and it is the only part of the certificate that the choice of procedure controls.

\begin{remark}[The functional at the parameters of Chapter~\ref{ch:certification}]\label{rem:sharp-func-numbers}
In the worked example of Chapter~\ref{ch:certification} the conformal part of the two-sided error at $n=25$, $\alpha=0.1$, $\delta=0.05$ is $\max\{0.9-0.8239,\,0.9856-0.9\}=0.0856$, set by the upper tail because the $\Beta(24,2)$ law is left-skewed and over-coverage is the more likely direction. The full two-sided bound of Proposition~\ref{prop:twosided}(c) is $0.1345$ once the drift $0.0391$ and the Markov term $0.0098$ are added. Definition~\ref{def:sharp-error} isolates the $0.0856$: that is the part the design of the procedure can move.
\end{remark}

\section{The optimal block length}\label{sec:sharp-block}

\begin{theorem}[Optimal block length for the blocked route]\label{thm:block-b}
Let the record of $N$ periods be blocked as in Algorithm~\ref{alg:prelim-blocking} of Chapter~\ref{ch:prelim}: blocks of length $b$ separated by gaps of length $b$, $n=\lfloor N/(2b)\rfloor$ blocks, threshold $\hat q=S_{(k)}$ with $k=\lceil(1-\alpha)(n+1)\rceil\le n$. Suppose Assumption~\ref{ass:mix} of Chapter~\ref{ch:certification} holds with $\beta(\cdot)\in\Bclass{r}(C)$, $r>1$, and that $F_P$ is continuous.
\begin{enumerate}[nosep,label=(\alph*)]
\item (The bound.) The expected calibration-conditional shortfall obeys
\begin{equation}\label{eq:sharp-shortfall}
\E\big[\big(1-\alpha-F_P(\hat q)\big)^+\big]\;\le\;\frac1{2\sqrt{n+2}}+(n+1)\beta(b)\;\le\;\sqrt{\frac b{2N}}+C\,N\,b^{-(r+1)}\qquad(2\le b\le N/2).
\end{equation}
\item (The minimiser.) Write $f(b):=\sqrt{b/N}+Nb^{-(r+1)}$ for the right side of \eqref{eq:sharp-shortfall} with the constants $1/\sqrt2$ and $C$ set to one. Then $f$ is strictly decreasing and then strictly increasing on $(0,\infty)$, and its unique minimiser and minimum value are
\begin{equation}\label{eq:sharp-bstar}
b^*=\big(2(r+1)\big)^{2/(2r+3)}N^{3/(2r+3)},\qquad
f(b^*)=\frac{2r+3}{2(r+1)}\big(2(r+1)\big)^{1/(2r+3)}N^{-r/(2r+3)} .
\end{equation}
\item (The balance point is not the minimiser.) The point $b_{\mathrm{bal}}:=N^{3/(2r+3)}$ at which the two terms of $f$ are equal satisfies $b^*/b_{\mathrm{bal}}=(2(r+1))^{2/(2r+3)}>1$ and $f(b_{\mathrm{bal}})>f(b^*)$.
\end{enumerate}
\end{theorem}

\begin{proof}
\emph{(a).} By Corollary~\ref{cor:coupled} of Chapter~\ref{ch:certification} there is an event $E_n$ of probability at least $1-n\beta(b)$ on which the calibration scores coincide with independent copies $S^*_{1:n}$ with law $P$ given the fitting sample, and on which $\hat q=\hat q^*$. On $E_n$, and conditionally on the fitting sample, $F_P(\hat q^*)$ has the $\Beta(k,n+1-k)$ law exactly, because $F_P$ is continuous (Theorem~\ref{thm:prelim-beta} of Chapter~\ref{ch:prelim}, via Corollary~\ref{cor:exact}(a) of Chapter~\ref{ch:certification}). Write $B$ for a $\Beta(k,n+1-k)$ variable. Since $k=\lceil(1-\alpha)(n+1)\rceil\ge(1-\alpha)(n+1)$, the mean of $B$ satisfies $\E B=k/(n+1)\ge1-\alpha$, hence
\[
\big(1-\alpha-B\big)^+\le\big(\E B-B\big)^+\le\big|B-\E B\big| ,
\]
and by the Cauchy--Schwarz inequality $\E|B-\E B|\le\operatorname{sd}(B)$. Lemma~\ref{lem:prelim-order-stat} of Chapter~\ref{ch:prelim} gives
\[
\Var B=\frac{k(n+1-k)}{(n+1)^2(n+2)}=\frac{\mu(1-\mu)}{n+2},\qquad \mu:=\frac k{n+1},
\]
and $\mu(1-\mu)\le\tfrac14$, so $\operatorname{sd}(B)\le1/(2\sqrt{n+2})$. On the complement of $E_n$ the integrand $(1-\alpha-F_P(\hat q))^+$ is at most $1-\alpha\le1$, contributing at most $\Prob(E_n^c)\le n\beta(b)$; adding the coupling of the deployment block, which Corollary~\ref{cor:coupled} charges at one further $\beta(b)$, gives the stated $(n+1)\beta(b)$ and the first inequality of \eqref{eq:sharp-shortfall}.

For the second inequality, $n=\lfloor N/(2b)\rfloor$ gives $n+2>N/(2b)$, so
\[
\frac1{2\sqrt{n+2}}<\frac12\sqrt{\frac{2b}N}=\sqrt{\frac b{2N}} .
\]
For the coupling term, $b\le N/2$ gives $N/(2b)\ge1$ and hence $n+1\le N/(2b)+1\le N/b$; with $\beta(b)\le Cb^{-r}$,
\[
(n+1)\beta(b)\le\frac Nb\cdot Cb^{-r}=C\,N\,b^{-(r+1)} .
\]

\emph{(b).} Write $f(b)=N^{-1/2}b^{1/2}+Nb^{-(r+1)}$, so that
\[
f'(b)=\tfrac12N^{-1/2}b^{-1/2}-(r+1)Nb^{-(r+2)},\qquad
b^{1/2}f'(b)=\tfrac12N^{-1/2}-(r+1)N\,b^{-(2r+3)/2} .
\]
The function $b\mapsto b^{1/2}f'(b)$ is strictly increasing on $(0,\infty)$, because $-b^{-(2r+3)/2}$ is, and it vanishes exactly where
\[
b^{(2r+3)/2}=2(r+1)N^{3/2}\iff b=\big(2(r+1)\big)^{2/(2r+3)}N^{3/(2r+3)}=b^* .
\]
Since $b^{1/2}>0$, the derivative $f'$ has the same sign as $b^{1/2}f'(b)$, so $f'<0$ on $(0,b^*)$ and $f'>0$ on $(b^*,\infty)$: the minimiser is unique and equal to $b^*$, and $f$ is strictly decreasing before it and strictly increasing after it. (Convexity is not needed and does not hold globally: $f''(b)=-\tfrac14N^{-1/2}b^{-3/2}+(r+1)(r+2)Nb^{-(r+3)}$ is positive only for $b^{(2r+3)/2}<4(r+1)(r+2)N^{3/2}$, an interval that contains $b^*$ but not the whole half-line.)

For the value, substitute. Write $\lambda:=2(r+1)$ and $\theta:=3/(2r+3)$, so $b^*=\lambda^{2/(2r+3)}N^\theta$. The first term is
\[
\sqrt{b^*/N}=\lambda^{1/(2r+3)}N^{(\theta-1)/2}=\lambda^{1/(2r+3)}N^{-r/(2r+3)} ,
\]
because $(\theta-1)/2=(3/(2r+3)-1)/2=-r/(2r+3)$. The second term is
\[
N(b^*)^{-(r+1)}=N\lambda^{-2(r+1)/(2r+3)}N^{-3(r+1)/(2r+3)}=\lambda^{-2(r+1)/(2r+3)}N^{-r/(2r+3)} ,
\]
because $1-3(r+1)/(2r+3)=(2r+3-3r-3)/(2r+3)=-r/(2r+3)$. Their ratio is
\[
\frac{N(b^*)^{-(r+1)}}{\sqrt{b^*/N}}=\lambda^{-(2r+3)/(2r+3)}=\lambda^{-1}=\frac1{2(r+1)} ,
\]
so at the minimiser the coupling term is exactly $1/(2(r+1))$ times the concentration term. Hence
\[
f(b^*)=\Big(1+\frac1{2(r+1)}\Big)\lambda^{1/(2r+3)}N^{-r/(2r+3)}=\frac{2r+3}{2(r+1)}\big(2(r+1)\big)^{1/(2r+3)}N^{-r/(2r+3)} ,
\]
which is \eqref{eq:sharp-bstar}.

\emph{(c).} At $b_{\mathrm{bal}}=N^{3/(2r+3)}$ the two terms of $f$ are equal: $\sqrt{b_{\mathrm{bal}}/N}=N^{(3/(2r+3)-1)/2}=N^{-r/(2r+3)}$ and $Nb_{\mathrm{bal}}^{-(r+1)}=N^{1-3(r+1)/(2r+3)}=N^{-r/(2r+3)}$, so $f(b_{\mathrm{bal}})=2N^{-r/(2r+3)}$. The ratio of minimiser to balance point is $b^*/b_{\mathrm{bal}}=\lambda^{2/(2r+3)}>1$ since $\lambda=2(r+1)>1$, and
\[
\frac{f(b_{\mathrm{bal}})}{f(b^*)}=\frac{2}{\frac{2r+3}{2(r+1)}\lambda^{1/(2r+3)}}=\frac{4(r+1)}{(2r+3)\,\big(2(r+1)\big)^{1/(2r+3)}}>1
\]
for every $r>1$, as the display in Example~\ref{ex:sharp-bstar} verifies numerically. The two terms of $f$ are not equally expensive at the optimum: the coupling term should be $2(r+1)$ times \emph{smaller} than the concentration term, because it falls in $b$ at the polynomial rate $r+1$ while the concentration term rises only as $\sqrt b$, and the minimiser of a sum sits where the derivatives cancel, not where the terms agree.
\end{proof}

\begin{example}[The minimiser at polynomial rate two and record length five hundred]\label{ex:sharp-bstar}
Take $r=2$, $N=500$, $C=1$. Then $2r+3=7$, $\lambda=6$, and
\[
b^*=6^{2/7}\times500^{3/7}=1.6685\times14.345=23.935,\qquad f(b^*)=\tfrac76\times6^{1/7}\times500^{-2/7}=1.5070\times0.16936=0.25526 .
\]
The balance point is $b_{\mathrm{bal}}=500^{3/7}=14.345$ with $f(b_{\mathrm{bal}})=2\times500^{-2/7}=0.33876$. The minimiser is larger than the balance point by the factor $1.67$, and the balance point costs $33$ per cent more than the minimum. Rounded to whole periods the recommendation is a block of $24$ periods and a gap of $24$, giving $n=\lfloor500/48\rfloor=10$ blocks.
\[
\begin{array}{lrrrr}
\toprule
r & 1.5 & 2 & 3 & 5\\
\midrule
b^* & 38.24 & 23.93 & 12.60 & 6.15\\
b_{\mathrm{bal}} & 22.36 & 14.35 & 7.94 & 4.20\\
b^*/b_{\mathrm{bal}} & 1.710 & 1.669 & 1.587 & 1.466\\
f(b^*) & 0.3318 & 0.2553 & 0.1786 & 0.1201\\
f(b_{\mathrm{bal}})/f(b^*) & 1.275 & 1.327 & 1.411 & 1.525\\
\bottomrule
\end{array}
\]
The block length falls as the mixing exponent rises, as it must: a process that forgets faster needs a shorter gap to decouple, and the record can afford more blocks. The penalty for using the balance point rises with $r$, because the coupling term then falls more steeply and the optimum moves further from equality.
\end{example}

\datanote{Every entry of the table in Example~\ref{ex:sharp-bstar} is the closed form \eqref{eq:sharp-bstar} evaluated in double precision at $N=500$; the two functions are implemented in the certificate package of Appendix~\ref{app:data} as \texttt{optimal\_block\_length} and \texttt{optimal\_deficit}, with the balance point kept under the separate name \texttt{balance\_block\_length} so that it cannot be read as optimal. The value of $f$ at the minimiser was checked against a direct grid search on $b\in[2,60]$ in steps of $10^{-3}$.}

\begin{remark}[The exact bound, with its constants, at the same parameters]\label{rem:sharp-exact-b}
The function $f$ discards the constants $1/\sqrt2$ and $C$ and replaces $n+2$ by $N/(2b)$. Minimising the first inequality of \eqref{eq:sharp-shortfall} instead, that is $g(b):=1/(2\sqrt{n+2})+(n+1)Cb^{-r}$ with $n=\lfloor N/(2b)\rfloor$, at $C=1$, $r=2$, $N=500$ gives an integer minimiser $b=25$ with $g(25)=0.1619$, against $g(24)=0.1634$, $g(23)=0.1651$, $g(20)=0.1661$ and $g(10)=0.3562$. The continuous relaxation $b^*=23.9$ therefore lands within one period of the integer minimiser of the exact bound, and the two disagree only through the floor in $n$, which makes $g$ piecewise constant in stretches. The practical rule is to evaluate $g$ on the integers near $b^*$ and take the best; the closed form \eqref{eq:sharp-bstar} is what tells one where to look and what the rate is.
\end{remark}

\begin{remark}[The rate, and what it is not]\label{rem:sharp-rate}
Theorem~\ref{thm:block-b} gives the blocked route an expected calibration-conditional shortfall of order $N^{-r/(2r+3)}$, with the explicit constant $\frac{2r+3}{2(r+1)}(2(r+1))^{1/(2r+3)}$, equal to $1.51$ at $r=2$. This rate is a property of the design, not of the problem. It is not minimax, and Section~\ref{sec:sharp-notminimax} proves that it is not; it is not the marginal rate, which \citet[Cor.~2]{barber2026} obtain as $N^{-r/(r+1)}$ without blocking; and it is not attained by the bare balance of the two terms, which Theorem~\ref{thm:block-b}(c) shows to be a third worse. The exponent $r/(2r+3)$ increases from $0.2$ at $r=1$ to $0.5$ as $r\to\infty$: no amount of mixing makes the blocked route better than $N^{-1/2}$, which is what an independent sample of $N$ scores would give without blocking at all. That observation is the content of the next section.
\end{remark}

\section{The blocked rate is not minimax}\label{sec:sharp-notminimax}

The blocked design pays for its exact Beta law with a factor $\sqrt{b}$ in effective sample size. This section proves that the payment is real and unnecessary as far as rate is concerned, by showing that the ordinary sample quantile of the whole unblocked record attains $(N\delta)^{-1/2}$ over every class $\Bclass{r}$ with $r>1$, and that no procedure can do better than $N^{-1/2}$ at fixed confidence.

\subsection{Variance of the empirical distribution function under mixing}

\begin{lemma}[Covariance of half-line indicators]\label{lem:sharp-cov}
Let $(S_t)$ be stationary with $\alpha$-mixing coefficients $\alpha(\cdot)$ and $\beta$-mixing coefficients $\beta(\cdot)$ in the sense of Definition~\ref{def:prelim-mixing} of Chapter~\ref{ch:prelim}. For every $x\in\R$, every $k\ge1$ and every choice of the closed or open half-line,
\[
\big|\Cov\big(\ind{S_0\le x},\ind{S_k\le x}\big)\big|\;\le\;\alpha(k)\;\le\;\beta(k).
\]
\end{lemma}

\begin{proof}
The covariance of two indicators is $\Prob(A\cap B)-\Prob(A)\Prob(B)$ with $A:=\{S_0\le x\}\in\sigma(S_0)$ and $B:=\{S_k\le x\}\in\sigma(S_k)$. Definition~\ref{def:prelim-ab-sigma} of Chapter~\ref{ch:prelim} defines $\alpha(\sigma(S_0),\sigma(S_k))$ as the supremum of exactly this quantity over $A$ and $B$ in the two $\sigma$-fields, so the covariance is at most $\alpha(\sigma(S_0),\sigma(S_k))\le\alpha(k)$ by Definition~\ref{def:prelim-mixing}. That $\alpha(k)\le\beta(k)$ is the last statement of Lemma~\ref{lem:prelim-beta-tv} of Chapter~\ref{ch:prelim}. The argument does not use the form of the sets, so it applies verbatim to $\{S_0<x\}$ and $\{S_k<x\}$.
\end{proof}

\begin{lemma}[Variance of the empirical distribution function]\label{lem:sharp-var}
Let $(S_t)$ be stationary in $\Bclass{r}(C)$ with $r>1$, and let $F_N(x):=N^{-1}\sum_{t=1}^N\ind{S_t\le x}$. Then for every $x\in\R$,
\[
\Var F_N(x)\;\le\;\frac{\tfrac14+2C\zeta(r)}N ,
\]
where $\zeta$ is the Riemann zeta function; the same bound holds for the left-limit version $F_N(x^-):=N^{-1}\sum_t\ind{S_t<x}$.
\end{lemma}

\begin{proof}
Expanding the variance and using stationarity,
\[
\Var F_N(x)=\frac1{N^2}\sum_{s,t=1}^N\Cov\big(\ind{S_s\le x},\ind{S_t\le x}\big)
=\frac1{N^2}\Big[N\Var\ind{S_1\le x}+2\sum_{k=1}^{N-1}(N-k)\Cov\big(\ind{S_0\le x},\ind{S_k\le x}\big)\Big].
\]
The variance of a Bernoulli indicator is at most $\tfrac14$. By Lemma~\ref{lem:sharp-cov} each covariance is at most $\beta(k)\le Ck^{-r}$ in absolute value, and $N-k\le N$, so
\[
\Var F_N(x)\le\frac1{N^2}\Big[\frac N4+2N\sum_{k=1}^{\infty}Ck^{-r}\Big]=\frac{\tfrac14+2C\zeta(r)}N ,
\]
the series converging because $r>1$. The left-limit version is identical by the last sentence of Lemma~\ref{lem:sharp-cov}.
\end{proof}

\begin{remark}[No blocking, and where the mixing class enters]\label{rem:sharp-var}
Lemma~\ref{lem:sharp-var} uses no coupling, no gaps and no blocks. The mixing class enters through one convergent series, and the effective sample size is $N$ up to the additive constant $2C\zeta(r)$, not $N/(2b)$. This is the whole reason the unblocked quantile beats the blocked one in rate: blocking discards half the record to gaps and then treats the remaining $n=\lfloor N/(2b)\rfloor$ blocks as the sample, whereas the covariance bound uses every observation and pays only a constant.
\end{remark}

\subsection{The upper bound for the unblocked sample quantile}

\begin{theorem}[The blocked rate is not minimax]\label{thm:block-rate}
Fix $\alpha\in(0,\tfrac12)$, $r>1$ and $C>0$, and let $\Bclass{r}(C)$ be as in Definition~\ref{def:sharp-error}. Let $T_N:=S_{(\lceil(1-\alpha)N\rceil)}$ be the sample quantile of the whole unblocked record.
\begin{enumerate}[nosep,label=(\alph*)]
\item (Upper bound.) For every stationary process in $\Bclass{r}(C)$ with continuous $F_P$ and every $\varepsilon>0$,
\begin{equation}\label{eq:sharp-cheb}
\Prob\big(\big|F_P(T_N)-(1-\alpha)\big|>\varepsilon\big)\;\le\;\frac{\tfrac12+4C\zeta(r)}{N\varepsilon^2},
\qquad\text{hence}\qquad
\mathcal E_\delta(T_N)\;\le\;\sqrt{\frac{\tfrac12+4C\zeta(r)}{N\delta}} .
\end{equation}
\item (Lower bound.) For every $\delta\in(0,\tfrac14)$ and every $N\ge4$,
\begin{equation}\label{eq:sharp-twopoint}
\mathcal E^*_\delta(N,r)\;\ge\;\frac{\sqrt2\,\alpha(1-4\delta)}{3}\,N^{-1/2} .
\end{equation}
\item (Consequence.) At fixed $\delta$ the minimax coverage error over $\Bclass{r}(C)$ is of order $N^{-1/2}$ for every $r>1$, attained without blocking; the blocked rate $N^{-r/(2r+3)}$ of Theorem~\ref{thm:block-b} is slower by the factor $\sqrt{b^*}$. Blocking buys the exact $\Beta(k,n+1-k)$ law and a certificate whose only dependence on $(C,r)$ is the failure probability $(n+1)\beta(b)$. It does not buy rate.
\end{enumerate}
\end{theorem}

\begin{proof}
\emph{(a).} Write $u:=1-\alpha$, $k_N:=\lceil uN\rceil$, and let $q_-:=\inf\{x:F_P(x)\ge u-\varepsilon\}$ and $q_+:=\inf\{x:F_P(x)\ge u+\varepsilon\}$, so that $F_P(q_\mp)=u\mp\varepsilon$ by continuity, assuming $0<u-\varepsilon$ and $u+\varepsilon<1$; if either fails, the event in \eqref{eq:sharp-cheb} is empty or the bound is weaker than the trivial one and there is nothing to prove.

Suppose first $F_P(T_N)<u-\varepsilon$. Then $T_N<q_-$, since $x\ge q_-$ implies $F_P(x)\ge u-\varepsilon$. At least $k_N$ of the $N$ scores are at most $T_N$, because $T_N$ is the $k_N$-th order statistic; hence
\[
F_N(q_-^-)\ \ge\ F_N(T_N)\ \ge\ \frac{k_N}N\ \ge\ u ,
\]
the first inequality because $T_N<q_-$ so $\{S_t\le T_N\}\subseteq\{S_t<q_-\}$. Since $F_P(q_-^-)=F_P(q_-)=u-\varepsilon$ by continuity, this forces
\[
F_N(q_-^-)-\E F_N(q_-^-)\ \ge\ \varepsilon .
\]
By Chebyshev's inequality and Lemma~\ref{lem:sharp-var} the probability of this is at most $(\tfrac14+2C\zeta(r))/(N\varepsilon^2)$.

Suppose next $F_P(T_N)>u+\varepsilon$. Then $T_N>q_+$: indeed $F_P(T_N)>u+\varepsilon=F_P(q_+)$ and $F_P$ is non-decreasing, so $T_N>q_+$. At most $k_N-1$ of the scores are strictly less than $T_N$, so
\[
F_N(q_+)\ \le\ \frac{k_N-1}N\ <\ u ,
\]
using $\{S_t\le q_+\}\subseteq\{S_t<T_N\}$ and $k_N=\lceil uN\rceil<uN+1$. Since $\E F_N(q_+)=F_P(q_+)=u+\varepsilon$, this forces $F_N(q_+)-\E F_N(q_+)<-\varepsilon$, again of probability at most $(\tfrac14+2C\zeta(r))/(N\varepsilon^2)$ by Chebyshev and Lemma~\ref{lem:sharp-var}.

Adding the two probabilities gives \eqref{eq:sharp-cheb}; setting the right side equal to $\delta$ and solving for $\varepsilon$ gives the stated bound on $\mathcal E_\delta(T_N)$.

\emph{(b).} Independent identically distributed scores have $\beta(k)=0$ for every $k\ge1$ and so belong to $\Bclass{r}(C)$ for every $C>0$ and $r>1$. It therefore suffices to bound the minimax error from below over independent samples, and a two-point construction does it.

Fix $\gamma\in(0,1/\sqrt2]$. Let $P_0$ be the uniform law on $[0,1]$ and let $P_1$ have density $1+\gamma$ on $[0,\tfrac12]$ and $1-\gamma$ on $(\tfrac12,1]$. Both have continuous distribution functions,
\[
F_0(t)=t,\qquad F_1(t)=\begin{cases}(1+\gamma)t,&t\in[0,\tfrac12],\\ t+\gamma(1-t),&t\in[\tfrac12,1],\end{cases}
\]
so that $\Delta(t):=F_1(t)-F_0(t)$ equals $\gamma t$ on $[0,\tfrac12]$ and $\gamma(1-t)$ on $[\tfrac12,1]$, and vanishes off $[0,1]$.

\emph{Separation.} Put $c:=\gamma\alpha/3$ and let $t\in\R\cup\{+\infty\}$ satisfy $|F_0(t)-u|<c$, where $u=1-\alpha\in(\tfrac12,1)$. For $t<0$ one has $|F_0(t)-u|=u\ge\tfrac12>c$, and for $t>1$ (including $t=+\infty$) one has $|F_0(t)-u|=\alpha>c$; so $t\in(u-c,u+c)\subseteq(\tfrac12,1)$ and $F_0(t)=t$. Then
\[
|F_1(t)-u|=\big|F_0(t)-u+\Delta(t)\big|\ \ge\ \gamma(1-t)-c\ >\ \gamma(1-u-c)-c=\gamma(\alpha-c)-c ,
\]
and $\gamma(\alpha-c)-c\ge c$ is equivalent to $\gamma\alpha\ge c(2+\gamma)$, which holds with $c=\gamma\alpha/3$ because $\gamma\le1$. Hence for every $t$,
\begin{equation}\label{eq:sharp-sep}
\max\big\{|F_0(t)-u|,\;|F_1(t)-u|\big\}\ \ge\ c=\frac{\gamma\alpha}3 .
\end{equation}

\emph{Reduction.} Let $T=T(S_{1:N})$ be any threshold procedure and set $A_j:=\{|F_j(T)-u|<c\}$ for $j=0,1$. Both are events in the sample space of the record, and by \eqref{eq:sharp-sep} they are disjoint. Writing $\Prob_j$ for the law of $N$ independent draws from $P_j$,
\[
\Prob_0\big(\mathcal E_0(T)\ge c\big)+\Prob_1\big(\mathcal E_1(T)\ge c\big)=\Prob_0(A_0^c)+\Prob_1(A_1^c)\ \ge\ \Prob_0(A_1)+\Prob_1(A_1^c)\ \ge\ 1-\dTV\big(\Prob_0,\Prob_1\big),
\]
the first inequality because $A_1\subseteq A_0^c$ and the second by the definition of total variation applied to the set $A_1$. Here $\mathcal E_j$ denotes the error functional of Definition~\ref{def:sharp-error} computed with $F_P=F_j$.

\emph{Distance.} The Kullback--Leibler divergence of a single observation is
\[
\mathrm{KL}(P_0\,\|\,P_1)=\int_0^{1/2}\log\frac1{1+\gamma}+\int_{1/2}^1\log\frac1{1-\gamma}=-\tfrac12\log(1-\gamma^2)\ \le\ \gamma^2 ,
\]
using $-\log(1-v)\le2v$ for $v\in[0,\tfrac12]$, valid since $\gamma^2\le\tfrac12$. Divergence is additive over independent coordinates, so $\mathrm{KL}(\Prob_0\|\Prob_1)\le N\gamma^2$, and Pinsker's inequality \citep[\S2.4.1]{tsybakov2009} gives $\dTV(\Prob_0,\Prob_1)\le\sqrt{N\gamma^2/2}=\gamma\sqrt{N/2}$.

\emph{Choice of $\gamma$.} Take $\gamma:=(1-4\delta)\sqrt{2/N}$, which lies in $(0,1/\sqrt2]$ as soon as $N\ge4$ because $1-4\delta<1$. Then $\dTV(\Prob_0,\Prob_1)\le1-4\delta$ and the reduction gives
\[
\max_{j\in\{0,1\}}\Prob_j\big(\mathcal E_j(T)\ge c\big)\ \ge\ \tfrac12\big(1-(1-4\delta)\big)=2\delta\ >\ \delta .
\]
For that $j$ and every $\varepsilon<c$ one has $\Prob_j(\mathcal E_j(T)>\varepsilon)\ge\Prob_j(\mathcal E_j(T)\ge c)>\delta$, so $\mathcal E_\delta(T)\ge c$ under $\Prob_j$. Since $T$ was arbitrary and both $\Prob_0$ and $\Prob_1$ lie in $\Bclass{r}(C)$,
\[
\mathcal E^*_\delta(N,r)\ \ge\ c=\frac{\gamma\alpha}3=\frac{\sqrt2\,\alpha(1-4\delta)}3N^{-1/2} ,
\]
which is \eqref{eq:sharp-twopoint}.

\emph{(c).} Parts (a) and (b) match in $N$ at fixed $\delta$, so $\mathcal E^*_\delta(N,r)\asymp N^{-1/2}$ with constants depending on $\alpha$, $\delta$, $C$ and $r$ but not on $N$, and the exponent does not involve $r$. Theorem~\ref{thm:block-b} gives the blocked route $N^{-r/(2r+3)}$, and
\[
\frac{N^{-r/(2r+3)}}{N^{-1/2}}=N^{(2r+3-2r)/(2(2r+3))}=N^{3/(2(2r+3))}=\sqrt{b_{\mathrm{bal}}}\asymp\sqrt{b^*} ,
\]
so the blocked error is larger by the square root of the block length, which is the number of periods each block consumes per score it produces, doubled by the gaps. The positive statements about blocking are those of Chapter~\ref{ch:certification}: the exact $\Beta(k,n+1-k)$ law of Corollary~\ref{cor:exact}, the drift charged on a fixed half-line by Lemma~\ref{lem:halfline}, and the single term $(n+1)\beta(b)$ through which the assumed class enters. None of them is a rate.
\end{proof}

\begin{remark}[Constants, and which route a bank should use at $N=500$]\label{rem:sharp-constants}
The two halves of Theorem~\ref{thm:block-rate} match in exponent and not in constant. At $\alpha=0.1$, $\delta=0.05$, $C=1$, $r=2$ the upper constant is $\sqrt{(\tfrac12+4\zeta(2))/\delta}=\sqrt{7.0797/0.05}=11.90$ and the lower constant is $\sqrt2\times0.1\times0.8/3=0.0377$: a factor of $316$ separates them, and closing it is not attempted here.

The gap is not academic. At $N=500$ the Chebyshev bound \eqref{eq:sharp-cheb} evaluates to $0.532$, which exceeds the largest possible value of $\mathcal E$ and is therefore vacuous, while the blocked bound of Remark~\ref{rem:sharp-exact-b} is $0.162$ at $b=25$. Evaluating both at a range of record lengths, with $C=1$, $r=2$, $\delta=0.05$:
\[
\begin{array}{lrrrrr}
\toprule
N & 500 & 10^4 & 10^5 & 4\times10^5 & 1.6\times10^6\\
\midrule
\text{blocked, }\min_bg(b) & 0.162 & 0.072 & 0.0377 & 0.0254 & 0.0171\\
\text{unblocked, \eqref{eq:sharp-cheb}} & 0.532 & 0.119 & 0.0376 & 0.0188 & 0.0094\\
\bottomrule
\end{array}
\]
The two bounds cross near $N=10^5$ periods, some four hundred years of daily data. Asymptotically the unblocked quantile wins and the win is real; at every record length a bank possesses, the blocked bound is the smaller of the two, because $\sqrt{b^*}$ grows like $N^{3/(2(2r+3))}$, slowly, while the constant $\sqrt{(\tfrac12+4C\zeta(r))/\delta}$ is paid immediately. The honest summary is that blocking is rate-suboptimal and constant-optimal in the range of interest, and that the choice between the two routes at a given $N$ is a numerical comparison of two proved bounds, not a comparison of exponents.
\end{remark}

\datanote{The table in Remark~\ref{rem:sharp-constants} was produced by evaluating $g(b)=1/(2\sqrt{n+2})+(n+1)Cb^{-r}$ with $n=\lfloor N/(2b)\rfloor$ over all integers $2\le b\le N/2$ and taking the minimum, and by evaluating $\sqrt{(\tfrac12+4C\zeta(r))/(N\delta)}$ in closed form, at $C=1$, $r=2$, $\delta=0.05$, $\zeta(2)=1.644934$. The minimising block lengths are $25$, $65$, $210$, $378$ and $692$ at the five record lengths. No data enter.}

\begin{remark}[What a matching lower bound over $\Bclass{r}$ would have to look like]\label{rem:sharp-whatnext}
Theorem~\ref{thm:block-rate}(b) is proved on independent scores, which lie in every class, so it says nothing about mixing: the exponent $r$ is invisible in it. That is not a defect of the proof but a property of the problem at fixed confidence. A lower bound in which $r$ appears must therefore be sought where the independent case is not the worst case, and Theorem~\ref{thm:block-rate}(a) says where that is: the Chebyshev bound degrades as $\delta^{-1/2}$, so a construction whose damage is rare and large, rather than common and small, can beat it at small $\delta$. Section~\ref{sec:sharp-lower} builds such a construction.
\end{remark}

\section{The mixing-driven lower bound at high confidence}\label{sec:sharp-lower}

\subsection{The statement}

\begin{theorem}[Mixing-driven lower bound at high confidence]\label{thm:block-lower}
Fix $r>1$ and $C>0$. There are constants $c_r>0$ and $\delta_0\in(0,\tfrac14)$, depending on $r$ and $C$ only, such that for every $N$ there is a strictly stationary process in $\Bclass{r}(C)$ with continuous marginal distribution function for which:
\begin{enumerate}[nosep,label=(\alph*)]
\item (Sample quantile.) the unblocked sample quantile $T_N$ of Theorem~\ref{thm:block-rate} satisfies, for every $\delta\in[N^{-r},\delta_0]$,
\[
\Prob\Big(\mathcal E(T_N)\ \ge\ c_r\,\big(N^r\delta\big)^{-1/(r+1)}\Big)\ \ge\ \delta;
\]
\item (All procedures.) the same holds for every threshold procedure $T=T(S_{1:N})$, so that
\[
\mathcal E^*_\delta(N,r)\ \ge\ c_r\,\big(N^r\delta\big)^{-1/(r+1)},\qquad \delta\in[N^{-r},\delta_0].
\]
\end{enumerate}
Together with an upper bound of Fuk--Nagaev type for bounded strongly mixing sequences applied to $F_N$ at the two quantiles $q_{1-\alpha\mp\varepsilon}(P)$, this gives
\begin{equation}\label{eq:sharp-minimax}
\mathcal E^*_\delta(N,r)\ \asymp\ \max\Big\{N^{-1/2}g_r(\delta),\ \big(N^r\delta\big)^{-1/(r+1)}\Big\},\qquad
\sqrt{\log(1/\delta)}\ \lesssim\ g_r(\delta)\ \lesssim\ \delta^{-1/r},
\end{equation}
and the mixing term binds exactly when $\delta<N^{-(r-1)/2}$.
\end{theorem}

Part (a) is a statement about one procedure and is provable by the route below. Part (b) is the doctoral claim and is not proved in this thesis; the display \eqref{eq:sharp-minimax} depends on both (b) and on the Fuk--Nagaev upper bound and inherits their status. The gap box at the end of this section says exactly what is missing.

\subsection{The construction}

\begin{definition}[Regime-renewal process]\label{def:sharp-renewal}
Fix two laws $P_1\ne P_2$ on $\R$ with continuous distribution functions and $\dK(P_1,P_2)=:\gamma_0>0$. Fix $r>1$ and a duration law on $\N$ with
\[
\Prob(D>m)=c_D\,m^{-(r+1)}\qquad(m\ge m_0),
\]
so that $\E D<\infty$. Let $(\theta_i)_{i\in\mathbb Z}$ be independent and uniform on $\{1,2\}$, and let $(D_i)_{i\in\mathbb Z}$ be independent copies of $D$, independent of the labels. Build a two-sided stationary renewal process with inter-arrival law $D$ by the standard delay construction: the interval covering the time origin is drawn from the size-biased law $\Prob(D=m)m/\E D$ and the position of the origin within it is uniform; the remaining intervals tile $\mathbb Z$ with independent copies of $D$. Label the $i$-th interval $\theta_i$. Conditionally on the renewal skeleton and the labels, the scores $(S_t)_{t\in\mathbb Z}$ are independent, with $S_t\sim P_{\theta_i}$ for every $t$ in the $i$-th interval.
\end{definition}

The process is strictly stationary, its marginal law is $P=\tfrac12P_1+\tfrac12P_2$, which is continuous, and it is a regenerative process whose regeneration epochs are the renewal epochs. Three properties matter.

\begin{lemma}[Mixing of the regime-renewal process]\label{lem:sharp-renewal-beta}
For the process of Definition~\ref{def:sharp-renewal} there is $\kappa\ge1$, not depending on $k$, with
\[
\beta(k)\ \le\ \kappa\,\Prob\big(\text{no renewal epoch in }(t,t+k]\big)\ =\ \kappa\,\Prob\big(R_t>k\big)
\ =\ \frac{\kappa}{\E D}\sum_{j>k}\Prob(D>j)\ \sim\ \frac{\kappa\,c_D}{r\,\E D}\,k^{-r},
\]
where $R_t$ is the residual life at $t$ under the stationary renewal measure. Hence the process lies in $\Bclass{r}(C)$ for every $C$ exceeding $\kappa c_D/(r\E D)$ and $N$ large enough.
\end{lemma}

\begin{route}
Every renewal epoch is a regeneration point: the label and duration of the interval that begins there are drawn afresh, independently of everything before. If a renewal occurs in $(t,t+k]$, couple the process after that epoch with an independent copy of the stationary process started from the same epoch; the two agree in law and, by the regeneration property, the coupled future is independent of $\sigma(S_s:s\le t)$. The disintegrated form of the $\beta$-coefficient (Lemma~\ref{lem:prelim-beta-tv} of Chapter~\ref{ch:prelim}) then bounds $\beta(k)$ by the probability that no renewal occurs in the gap, up to a factor accounting for the fact that the epoch itself is measurable with respect to the past through the residual life at $t$. The residual-life computation is the elementary identity $\Prob(R_t>k)=(\E D)^{-1}\sum_{j>k}\Prob(D>j)$ for the stationary renewal measure, and the tail asymptotics follow from $\Prob(D>j)=c_Dj^{-(r+1)}$ by comparison with $\int_k^\infty x^{-(r+1)}dx=k^{-r}/r$. The same device, sharing the refresh variables between two copies and identifying the disagreement event, is used in the proof of Proposition~\ref{prop:noest} of Chapter~\ref{ch:certification}.
\end{route}

\begin{remark}[The damage a long run does]\label{rem:sharp-run}
Suppose the record $S_{1:N}$ contains an interval of length $m$ carrying a single label, say $2$. Conditionally on that, the empirical distribution function is a mixture,
\[
F_N\ \approx\ \Big(1-\frac mN\Big)F_P+\frac mN F_2 = F_P+\frac mN\big(F_2-F_P\big),
\]
so at any fixed point $F_N$ is displaced from $F_P$ by up to $(m/N)\dK(P_2,P)=(m/N)\gamma_0/2$. The sample quantile $T_N$ solves $F_N(T_N)\approx1-\alpha$, so $F_P(T_N)\approx1-\alpha-(m/N)(F_2-F_P)(q_{1-\alpha})$ and the coverage error is of order $m/N$. A run is therefore worth its length as a fraction of the record, and nothing more; what makes it dangerous is that a heavy-tailed duration law makes long runs far more likely than the Gaussian fluctuation of $F_N$ would suggest.
\end{remark}

\begin{remark}[The trade-off that produces the rate]\label{rem:sharp-tradeoff}
The record of $N$ periods contains about $N/\E D$ intervals. The probability that at least one of them has length exceeding $m$ is of order
\[
\frac N{\E D}\Prob(D>m)\ \asymp\ N\,m^{-(r+1)} .
\]
Setting this equal to $\delta$ gives the critical run length
\[
m_\delta\ \asymp\ \Big(\frac N\delta\Big)^{1/(r+1)} ,
\]
and Remark~\ref{rem:sharp-run} then gives a coverage error of order
\[
\frac{m_\delta}N\ \asymp\ N^{1/(r+1)-1}\delta^{-1/(r+1)}=N^{-r/(r+1)}\delta^{-1/(r+1)}=\big(N^r\delta\big)^{-1/(r+1)} ,
\]
which is the rate of Theorem~\ref{thm:block-lower}. At $N=500$, $r=2$, $\delta=0.05$ the critical run is $m_\delta=(500/0.05)^{1/3}=21.5$ periods and the error is $21.5/500=0.043$, against the Gaussian scale $N^{-1/2}=0.045$: the two are of the same size at exactly this confidence, as the crossover condition below requires.
\end{remark}

\subsection{Proof route}

\begin{route}
\emph{Part (a), the sample quantile.} Fix $\delta\in[N^{-r},\delta_0]$ and set $m:=\lceil(N/\delta)^{1/(r+1)}\rceil$. Let $A$ be the event that some interval of the renewal skeleton lies entirely inside $\{1,\dots,N\}$ and has length at least $m$. By the renewal computation of Remark~\ref{rem:sharp-tradeoff}, $\Prob(A)\ge\delta$ for a suitable choice of $c_D$ and $\delta_0$; the restriction $\delta\ge N^{-r}$ keeps $m\le N$, so the event is not empty. On $A$, condition on the position and label of the long interval. The remaining $N-m$ scores are those of a stationary regime-renewal process with the long interval excised, and $F_N$ restricted to them concentrates at scale $(N-m)^{-1/2}\ll m/N$ by Lemma~\ref{lem:sharp-var} applied to the excised record, because $m/N\asymp(N^r\delta)^{-1/(r+1)}$ exceeds $N^{-1/2}$ throughout the range $\delta<N^{-(r-1)/2}$, which contains $[N^{-r},\delta_0]$ for $\delta_0$ small. Hence, on $A$ and off an event of small conditional probability, the displacement of Remark~\ref{rem:sharp-run} is not cancelled, and $\mathcal E(T_N)\ge c_r m/N$ with $c_r$ a fixed multiple of $\gamma_0$ times the density of $P$ near its $(1-\alpha)$-quantile. Choosing $P_1,P_2$ so that $F_2-F_P$ is bounded away from zero at $q_{1-\alpha}(P)$ fixes $c_r$.

\emph{Part (b), all procedures.} The step from one procedure to all procedures is a two-point argument in the style of Theorem~\ref{thm:block-rate}(b), with the two hypotheses being two regime-renewal processes rather than two independent laws. Take $\Prob^{(0)}$ to be the process of Definition~\ref{def:sharp-renewal} with the duration law truncated at $m$, and $\Prob^{(1)}$ the untruncated process. The two agree on every record containing no run longer than $m$, so
\[
\dTV\big(\Prob^{(0)}_{1:N},\Prob^{(1)}_{1:N}\big)\ \le\ \Prob^{(1)}(A)\ \asymp\ N\,m^{-(r+1)}\ \asymp\ \delta .
\]
Under $\Prob^{(0)}$ no run exceeds $m$, the empirical distribution function concentrates at the Gaussian scale, and a procedure tuned to $\Prob^{(0)}$ can achieve error $O(N^{-1/2})$. Under $\Prob^{(1)}$ a run of length $m$ occurs with probability of order $\delta$, and by Remark~\ref{rem:sharp-run} it displaces the coverage of any threshold that was accurate under $\Prob^{(0)}$ by $\asymp m/N$. The reduction of Theorem~\ref{thm:block-rate}(b), or the coverage-reduction lemma \citep[Lemma~5]{halkiewicz2026} in the form that converts a two-point testing bound into a coverage bound, then yields the claim provided the separation survives the conditioning: what must be shown is that no procedure measurable with respect to the record can detect the long run reliably enough to correct for it, which is exactly the statement that the two $N$-sample laws are $\delta$-close while the required correction is $\asymp m/N$.

\emph{The upper bound and the crossover.} For the matching upper bound, apply the Fuk--Nagaev inequality for bounded strongly mixing sequences \citep{rio2017} to the centred sums $\sum_t(\ind{S_t\le q}-F_P(q))$ at $q=q_{1-\alpha-\varepsilon}(P)$ and $q=q_{1-\alpha+\varepsilon}(P)$, exactly as Chebyshev was applied at those two points in the proof of Theorem~\ref{thm:block-rate}(a). The inequality has the shape of a Gaussian-type term $(1+N\varepsilon^2/(r\sigma^2))^{-r/2}$ plus a polynomial term $N^{-r}\varepsilon^{-(r+1)}$. Setting the polynomial term equal to $\delta$ reproduces $\varepsilon\asymp(N^r\delta)^{-1/(r+1)}$, matching the lower bound; setting the Gaussian-type term equal to $\delta$ gives $\varepsilon\asymp N^{-1/2}g_r(\delta)$ with $g_r$ trapped between $\sqrt{\log(1/\delta)}$, the sub-Gaussian value, and $\delta^{-1/r}$, the value the stated form of the inequality delivers. The two terms are equal when $(N^r\delta)^{-1/(r+1)}=N^{-1/2}$, that is when $\delta=N^{-(r-1)/2}$; the mixing term dominates for smaller $\delta$. At $N=500$, $r=2$ the threshold is $\delta=500^{-1/2}=0.045$, so every certificate written at the conventional confidence levels of $5$ per cent or better sits in the regime where the mixing exponent, and not the Gaussian scale, sets the achievable coverage error.
\end{route}

\begin{remark}[Consistency check against a published rate]\label{rem:sharp-bp-check}
A high-confidence lower bound must be consistent with the expected error, and this one is. Integrating the claimed rate over the confidence level,
\[
\E\,\mathcal E(T)\ \gtrsim\ \int_0^1\big(N^r\delta\big)^{-1/(r+1)}\,d\delta
= N^{-r/(r+1)}\int_0^1\delta^{-1/(r+1)}\,d\delta
= \frac{r+1}r\,N^{-r/(r+1)} ,
\]
the integral converging because $1/(r+1)<1$. The rate $N^{-r/(r+1)}$ is exactly what \citet[Cor.~1, eq.~(6)]{barber2026} obtain for the marginal coverage deficit of unblocked split conformal over $\Bclass{r}$, by minimising $\tau/N+2\beta(\tau)$ over the free lag $\tau$, which balances at $\tau\asymp N^{1/(r+1)}$. Two independent routes therefore produce the same exponent for the expected error, one from above for a particular procedure and one from below for all procedures; the disagreement between them is confined to the confidence level at which the error is measured, which is where Theorem~\ref{thm:block-lower} does its work. This is the strongest evidence available, short of the proof, that the construction of Definition~\ref{def:sharp-renewal} is the right one.
\end{remark}

\begin{remark}[Why this is the honest novelty claim]\label{rem:sharp-open}
\citet[Sec.~4]{barber2026} close their paper by asking ``whether split conformal is optimal for this class of problems, or whether alternative methods (that may use knowledge of the mixing properties of the time series) could improve these bounds''. Their own Theorem~2 is a construction for split conformal itself, tight up to a factor $(1-\alpha)/4$, and not a bound over procedures. Theorem~\ref{thm:block-lower}(b) is the calibration-conditional analogue of the question they leave open, and answering it is the claim of this chapter. It is not a claim to have produced the first optimal-design result for conformal prediction under dependence: \citet[Thm.~4]{halkiewicz2026} proves a minimax lower bound of exactly this shape for a rolling-origin procedure, and Section~\ref{sec:sharp-positioning} sets out the difference, which is that his hardness is drift-driven and this one is mixing-driven.
\end{remark}

\gap{Not established in this thesis. (1) Lemma~\ref{lem:sharp-renewal-beta}: the constant $\kappa$ in $\beta(k)\le\kappa\Prob(R_t>k)$ is not computed. The regeneration epoch after time $t$ is measurable with respect to the past through the residual life, so the naive identification of the $\beta$-coefficient with the no-renewal probability is not immediate and may require coupling two stationary renewal processes, with a constant that has not been extracted. (2) Theorem~\ref{thm:block-lower}(a): the concentration of the excised record, and hence the claim that the displacement of Remark~\ref{rem:sharp-run} is not cancelled, is asserted at the level of orders and not proved; the conditioning on the position and label of the long interval destroys stationarity of the remainder and Lemma~\ref{lem:sharp-var} does not apply to it verbatim. The constant $c_r$ is not computed and $\delta_0$ is not exhibited. (3) Theorem~\ref{thm:block-lower}(b): the two-point reduction over all procedures is a route only. It may lose a logarithmic factor, or hold only over the regime-renewal subclass of $\Bclass{r}$ rather than over the whole class; the step that must be supplied is that no record-measurable procedure can detect a run of length $m_\delta$ well enough to correct for it, given that the two $N$-sample laws are $\delta$-close. (4) The $\delta$-dependence $g_r$ in \eqref{eq:sharp-minimax} is open: it is trapped between $\sqrt{\log(1/\delta)}$ and $\delta^{-1/r}$ and its true order is not determined. (5) No verified locator is available for the Fuk--Nagaev inequality: the form quoted in the route is the shape of the inequality for bounded strongly mixing sequences as recorded in \cite{rio2017}, and no theorem or page number of that source has been checked against the printed text, so the upper half of \eqref{eq:sharp-minimax} rests on an unverified citation as well as on an unwritten proof.}

\section{Partially observed drift}\label{sec:sharp-drift}

Theorem~\ref{thm:sharp} settles the price of drift when the deployment law is unknown and the bank assumes a bound. In use the bank eventually sees deployment scores, and the question changes: how small a drift charge can be \emph{certified} from $m$ observed deployment scores, and can that be done without knowing the true drift in advance?

\begin{theorem}[Minimax certifiable deficit under partially observed drift]\label{thm:block-drift}
Let $P$ be the calibration score law, let $Q$ be the deployment score law with $\dK(P,Q)\le\varepsilon$ for an unknown $\varepsilon$, and suppose $m$ deployment scores $S^{\mathrm{dep}}_{1:m}$ are observed in addition to the calibration record. Call a pair $(T,\hat L)$, with $T$ a threshold and $\hat L=\hat L(S_{1:N},S^{\mathrm{dep}}_{1:m})$ a claimed coverage level, \emph{valid at confidence $1-\delta$} if
\[
\Prob\big(F_Q(T)\ \ge\ \hat L\big)\ \ge\ 1-\delta\qquad\text{for every }Q\text{ with }\dK(P,Q)\le\varepsilon\text{ and every }\varepsilon>0 .
\]
Then the largest achievable claimed level is
\[
\sup\big\{\E\hat L:(T,\hat L)\text{ valid at confidence }1-\delta\big\}
\;=\;(1-\alpha)-\varepsilon-\Theta\Big(\sqrt{\tfrac{\log(1/\delta)}m}\Big)
\]
up to the terms already charged in Theorem~\ref{thm:cov}(ii), and the supremum is attained adaptively, that is by a certificate that does not use $\varepsilon$.
\end{theorem}

\begin{route}
The statement has three parts and each has a route.

\emph{The $\varepsilon$ term is unimprovable.} This is Theorem~\ref{thm:sharp}: for the transported law $Q_M$ of Definition~\ref{def:sharp-transport} the coverage is lower than $C_T(P)$ by $\varepsilon$ up to corrections that vanish, so no valid certificate can claim more than $(1-\alpha)-\varepsilon$ even if $\varepsilon$ is known exactly.

\emph{The statistical term is a two-point bound on the deployment sample.} Here, unlike in Theorem~\ref{thm:sharp}, the adversary must choose a $Q$ that is both damaging and hard to distinguish from a benign one on the basis of $m$ observations. Take $Q_0$ with $\dK(P,Q_0)=\varepsilon$ and $Q_1$ with $\dK(P,Q_1)=\varepsilon+\eta$, the two differing only in a region of the line where they place mass of order $\eta$, so that $\mathrm{KL}(Q_0\|Q_1)\asymp\eta^2$ and $\dTV(Q_0^{\otimes m},Q_1^{\otimes m})\lesssim\eta\sqrt m$. A certificate valid under both must charge for $\varepsilon+\eta$ whenever $\eta\lesssim\sqrt{\log(1/\delta)/m}$, which is the claimed statistical term; the reduction is Le Cam's two-point method \citep[\S2.3, pp.~81--83]{tsybakov2009} in the form used in the proof of Theorem~\ref{thm:block-rate}(b).

\emph{Attainment, and adaptivity.} The matching certificate estimates the drift directly rather than through a transport distance. The two-sample Kolmogorov statistic $\hat\varepsilon_{\mathrm K}:=\sup_x|F_{P_n}(x)-F_{Q_m}(x)|$ has expectation $\dK(P,Q)$ up to the two empirical errors, and under independent sampling the \gls{dkw} bounds each at $\sqrt{\log(2/\delta)/(2m)}$. The certificate that reports $b_{n,k}(\delta)-\hat\varepsilon_{\mathrm K}-\sqrt{\log(2/\delta)/(2m)}-\beta(b)/\delta'$ never uses $\varepsilon$ and so is adaptive by construction. Its rate in $m$ is $m^{-1/2}$, against the $m^{-1/4}$ of the Wasserstein plug-in of Corollary~\ref{cor:plugin} of Chapter~\ref{ch:certification}: the square root in Proposition~\ref{prop:w1} is what costs the exponent, and it is avoided by estimating the Kolmogorov distance itself when a deployment sample exists. The Wasserstein route remains the only one available before deployment, when no sample exists and the drift must be assumed.

\emph{Where the mixing class re-enters.} The deployment scores are not independent. They come from consecutive or overlapping blocks of the deployment window, so the \gls{dkw} step must be replaced by a dependent-data bound: either blocking the deployment sample into $\lfloor m/(2b)\rfloor$ blocks, which costs the familiar $\sqrt b$ and a coupling term, or a direct uniform bound under $\Bclass{r}$.
\end{route}

\gap{Not established. (1) The upper bound: no matching valid certificate is constructed and no constant is computed; the display in Theorem~\ref{thm:block-drift} is written with $\Theta$ and both halves are routes. (2) The two-point construction is described but not exhibited: the pair $(Q_0,Q_1)$ must be built so that the Kolmogorov distances differ by exactly $\eta$ while the $m$-sample laws stay $\delta$-close, and the interaction with the density bound $L$ that Proposition~\ref{prop:w1} of Chapter~\ref{ch:certification} needs is not analysed. (3) The dependent replacement for the Dvoretzky--Kiefer--Wolfowitz step is not supplied, so the rate $m^{-1/2}$ is claimed only for independent deployment scores, which is not the case in the application. (4) Whether the adaptive certificate loses a factor $\sqrt{\log(1/\delta)}$ relative to the lower bound, or matches it, is not determined. (5) The statement quantifies over all $Q$ in a $\dK$-ball around a \emph{known} $P$; in the application $P$ is itself estimated from the calibration record, and the error in $P$ is not carried through.}

\section{The dependent empirical Wasserstein rate}\label{sec:sharp-w1}

Chapter~\ref{ch:certification} left one gap that is closed by citation rather than by proof. Corollary~\ref{cor:plugin} of that chapter certifies the calibration side of the Wasserstein plug-in, using the one-dimensional rate $\E\Wone(P_n,P)\le J_1(P)/\sqrt n$ for independent samples (Theorem~\ref{thm:prelim-bobkov} of Chapter~\ref{ch:prelim}), and leaves the deployment side uncontrolled because the deployment scores are dependent and not coupled to independent copies. The dependent analogue is in print.

\begin{proposition}[Dependent empirical Wasserstein rate over $\Bclass{r}$; specialised, not new]\label{prop:block-w1}
Let $(S_t)$ be a strictly stationary sequence with marginal law $Q$ supported in $[-M,M]$, let $Q_m$ be the empirical law of $S_{1:m}$, and suppose the block-score process has $\beta$-mixing coefficients $\beta_S(k)\le C'k^{-r}$ for $k\ge k_1$ and $\beta_S(k)\le1$ otherwise, with $r>1$. Then
\begin{equation}\label{eq:sharp-w1}
\big\|\Wone(Q_m,Q)\big\|_2\ \le\ \frac{2\sqrt2\,M}{\sqrt m}\Big(\sqrt{k_1}+\sqrt{\kappa\,C'\zeta(r)}\Big),
\end{equation}
where $\kappa\ge1$ is the absolute constant relating the weak-dependence coefficients of \citet[\S1]{dedecker2017} to the strong mixing coefficient. In particular $\E\Wone(Q_m,Q)\le c_r m^{-1/2}$ with $c_r$ explicit in $(M,C',r,k_1)$.
\end{proposition}

\begin{proof}
The bound is a specialisation of \citet[Prop.~3.2, eqs.~(3.7)--(3.8)]{dedecker2017}, whose second display reads
\[
\big\|\Wone(\mu_m,\mu)\big\|_2\ \le\ \frac{2\sqrt2}{\sqrt m}\int_0^\infty\sqrt{S_{\alpha,m}(t)}\,dt,\qquad
S_{\alpha,m}(t):=\sum_{k=0}^m\min\big\{\alpha_{1,X}(k),\,H(t)\big\},\quad H(t):=\Prob(|X_0|>t),
\]
for their weak-dependence coefficients $\alpha_{1,X}(k)$, which are built from indicators of half-lines and are dominated by the strong mixing coefficient up to the absolute constant $\kappa$: their class contains the strongly mixing sequences in the sense of Rosenblatt \citep[\S1]{dedecker2017}, and $\alpha(k)\le\beta(k)$ by Lemma~\ref{lem:prelim-beta-tv} of Chapter~\ref{ch:prelim}, so $\alpha_{1,X}(k)\le\kappa\beta_S(k)$.

Split the sum defining $S_{\alpha,m}$. For $k<k_1$ use $\min\{\alpha_{1,X}(k),H(t)\}\le H(t)$, giving at most $k_1H(t)$ in total. For $k\ge k_1$ use $\alpha_{1,X}(k)\le\kappa C'k^{-r}$ and sum, giving at most $\kappa C'\zeta(r)$. Hence $S_{\alpha,m}(t)\le k_1H(t)+\kappa C'\zeta(r)$ and, by subadditivity of the square root,
\[
\sqrt{S_{\alpha,m}(t)}\ \le\ \sqrt{k_1}\sqrt{H(t)}+\sqrt{\kappa C'\zeta(r)} .
\]
Because the scores are supported in $[-M,M]$, $H(t)=0$ for $t>M$, so the integral runs over $[0,M]$ only and $\sqrt{H(t)}\le1$ there. Therefore $\int_0^\infty\sqrt{S_{\alpha,m}(t)}\,dt\le M(\sqrt{k_1}+\sqrt{\kappa C'\zeta(r)})$, which gives \eqref{eq:sharp-w1}. The bound on the expectation follows from $\E X\le\|X\|_2$.
\end{proof}

\begin{remark}[Applying it to the deployment blocks]\label{rem:sharp-w1-apply}
The deployment scores of Chapter~\ref{ch:certification} are block scores computed from overlapping windows of length $b+1$ of the state process. By Lemma~\ref{lem:prelim-inherit} of Chapter~\ref{ch:prelim} the block-score process inherits $\beta_S(k)\le\beta(k-b)$ for $k>b$, so if the state process lies in $\Bclass{r}(C)$ then $\beta_S(k)\le C(k-b)^{-r}\le C2^rk^{-r}$ for $k\ge2b$. Proposition~\ref{prop:block-w1} therefore applies with $k_1=2b$ and $C'=C2^r$, and \eqref{eq:sharp-w1} becomes
\[
\big\|\Wone(Q_m,Q)\big\|_2\ \le\ \frac{2\sqrt2\,M}{\sqrt m}\Big(\sqrt{2b}+2^{r/2}\sqrt{\kappa\,C\zeta(r)}\Big) .
\]
The block length appears as $\sqrt b$ once more, and for the same reason as everywhere else in this chapter: a block of length $b$ needs a gap of length $b$ before the dependence has decayed, and the effective number of independent observations is reduced accordingly.
\end{remark}

\begin{example}[Size of the constant]\label{ex:sharp-w1-size}
Take scores normalised to $[-1,1]$, so $M=1$, with $b=10$, $C=1$, $r=2$ and $\kappa=1$. Then $\sqrt{2b}=4.472$ and $2^{r/2}\sqrt{\zeta(2)}=2\times1.283=2.565$, so the bracket is $7.037$ and
\[
\big\|\Wone(Q_m,Q)\big\|_2\ \le\ \frac{2\sqrt2\times7.037}{\sqrt m}=\frac{19.90}{\sqrt m},
\]
which is $0.890$ at $m=500$ and $0.089$ at $m=5\times10^4$. The bound is therefore of order one on a record of the length a converted book will have in its first two years, and becomes informative only at $m$ of order $10^4$. With $\kappa=4$ the constant rises to $27.2$ and the useful record length to about $7\times10^4$.
\end{example}

\datanote{The two constants in Example~\ref{ex:sharp-w1-size} are $2\sqrt2(\sqrt{20}+\sqrt{4\times1.644934})=19.90$ and $2\sqrt2(\sqrt{20}+\sqrt{16\times1.644934})=27.16$, evaluated in double precision with $\zeta(2)=1.644934$; the record lengths solve $19.90/\sqrt m=0.1$ and $27.16/\sqrt m=0.1$, giving $m=3.96\times10^4$ and $7.38\times10^4$. No data enter.}

\begin{remark}[Attribution]\label{rem:sharp-w1-attrib}
Nothing in Proposition~\ref{prop:block-w1} is new. The bound is \citep[Prop.~3.2, eq.~(3.8)]{dedecker2017}; the $m^{-1/2}$ rate under summable dependence is their Remark~3.4(1) \citep[Remark~3.4(1)]{dedecker2017}, and the polynomial case that matches $\Bclass{r}$ with $a=r$ is their Remark~3.4(2) \citep[Remark~3.4(2), eq.~(3.12)]{dedecker2017}, from which the same rate follows for bounded scores by the quantile-integral route, since $(a+1)/(2a)<1$ for $a>1$ makes their second integral converge. What this section contributes is the specialisation: the choice $a=r$, the reduction of the quantile integrals to the bracket of \eqref{eq:sharp-w1} for bounded scores, and the accounting of the block length through Lemma~\ref{lem:prelim-inherit}. The result closes the gap recorded in Section~\ref{sec:cert-drift} of Chapter~\ref{ch:certification} up to the constant $\kappa$.
\end{remark}

\gap{Not established: the absolute constant $\kappa$ relating the weak-dependence coefficients $\alpha_{1,X}(k)$ of \cite{dedecker2017} to the strong mixing coefficient $\alpha(k)$. Their definition is built from indicators of half-lines and the constant follows from it, but it has not been read off the source, so $\kappa$ appears as a symbol in \eqref{eq:sharp-w1} and Example~\ref{ex:sharp-w1-size} reports the bound for $\kappa=1$ and $\kappa=4$. Also not established: whether the constant in \eqref{eq:sharp-w1} can be improved for the specific case of overlapping block scores, where the dependence within a window of length $b$ is complete rather than merely bounded by one.}

\section{Positioning}\label{sec:sharp-positioning}

This section says, for each of the five nearest papers, what it has and what this chapter adds. No claim of priority is made for the template, for the object, or for the use of a lower bound.

\paragraph{\citet{oliveira2024}.} The blocking-and-coupling route for split conformal prediction on $\beta$-mixing data is theirs, built on Berbee's lemma and Lemma~4.1 of \citet{yu1994}. Their predictor is unblocked; blocking lives inside the proof \citep[Prop.~10, pp.~20--21]{oliveira2024}, and they say so explicitly \citep[p.~2]{oliveira2024}. They do optimise block sizes, taking $m=n^\lambda$ and $a=n^{1-\lambda}/2$ inside the bound, and they are careful to note that ``this optimization of block sizes plays a role exclusively on the coverage guarantees'' \citep[p.~8]{oliveira2024}. What is added here is the object being optimised and the exactness of the answer: their optimisation is of the marginal deficit, over exponents; Theorem~\ref{thm:block-b} optimises the expected calibration-conditional shortfall implied by the Beta law, over block lengths, and produces the exact minimiser with its constant and the statement that the balance point is not it. They prove no lower bound.

\paragraph{\citet{barber2026}.} For marginal coverage their bound dominates the blocked route: the deficit of the unblocked split-conformal threshold with a trained score on stationary $\beta$-mixing data is $2\beta(\tau)+2\beta(\tau_*)$ plus a term linear in the lags over the sample size \citep[Cor.~2]{barber2026}, of order $N^{-r/(r+1)}$ over $\Bclass{r}$, and they remark that their bound scales linearly in $\tau/n$ where the blocked bound scales as its square root \citep[Sec.~3.4]{barber2026}. Their Theorem~2 shows their guarantee tight for split conformal up to a factor $(1-\alpha)/4$, as a construction for that procedure and not as a bound over procedures. All their results are marginal; the paper contains no calibration-conditional statement. What is added here is the calibration-conditional object of Definition~\ref{def:sharp-error}, the honest accounting of Theorem~\ref{thm:block-rate} that says the blocked rate is not minimax in that object either, and the attempt in Theorem~\ref{thm:block-lower} on the question they leave open \citep[Sec.~4]{barber2026}. Remark~\ref{rem:sharp-bp-check} checks the claimed high-confidence rate against their expected-error rate and finds the two consistent.

\paragraph{\citet{halkiewicz2026}.} The template of this chapter is his: a coverage-error decomposition, an optimal design parameter obtained by balancing a noise term against a bias term, and a Le Cam two-point lower bound showing the resulting rate minimax over a class. He obtains an optimal calibration window $m^*\asymp T^{2\beta/(2\beta+1)}$ with coverage-error rate $T^{-\beta/(2\beta+1)}$ \citep[Sect.~3.4, eq.~(6)]{halkiewicz2026} and a matching minimax lower bound over every conformal prediction procedure \citep[Thm.~4]{halkiewicz2026}. Three differences are substantive and none is a matter of technique. His hardness is drift-driven: the lower bound is produced by a Gaussian mean bump over the last observations, and the class $\mathcal F(L,\beta)$ is indexed by the smoothness of the drift, with mixing entering only the upper bound. The hardness of Theorem~\ref{thm:block-lower} is mixing-driven: the class is indexed by the mixing exponent, the drift is charged separately and exactly by Theorem~\ref{thm:sharp}, and the construction of Definition~\ref{def:sharp-renewal} is stationary. His dependence notion is $\alpha$-mixing on the score process and he uses no coupling; the certificate of Chapter~\ref{ch:certification} needs $\beta$-mixing because the coupling needs total variation. And his object is the marginal coverage error, with the guarantee approximate rather than finite-sample \citep[Remark~3]{halkiewicz2026}, where Definition~\ref{def:sharp-error} takes the calibration-conditional error. This chapter is therefore not the first optimal-design result for conformal prediction under dependence and must not be read as one; its claim is the mixing-indexed lower bound for the conditional object.

\paragraph{\citet{ramos2026}.} The mechanism behind the drift term is in print as their Proposition~13: with the test score decoupled from the calibration block the realised coverage is $F_Q\circ F_P^{-1}(U_{(k)})$, so that $\sup_u|h(u)-u|=\dK(P,Q)$ is exactly the penalty \citep[Prop.~13, p.~11]{ramos2026}, and they name the blocking of \citet{yu1994} as one way to decouple. Their Proposition~15 reaches the Beta law asymptotically under $\alpha$-mixing through a Berry--Esseen bound; Theorem~\ref{thm:cov}(ii) of Chapter~\ref{ch:certification} reaches it non-asymptotically under $\beta$-mixing through the coupling, which is the difference between the two routes. Their Theorem~1 supplies the exchangeable rank bracket used in the proof of Proposition~\ref{prop:twosided}(a) of Chapter~\ref{ch:certification}. What is added here is Theorem~\ref{thm:sharp}: the attained constant $1$ by transport of the bottom mass to a far point is not in their paper, and the transport mechanism they describe is the map, not the extremal choice of $Q$.

\paragraph{\citet{zheng2024}.} Theirs is the closest Markov-chain analogue of blocking: $K$-split thinning of the calibration set with a gap of the order of the mixing time, coverage gap reduced from $t_{\mathrm{mix}}\log n/n$ to $t_{\mathrm{mix}}/(n\log n)$, with the cost of thinning being the number of retained points times $\beta(K)$ \citep[Prop.~5.1, p.~61475]{zheng2024}. They do optimise their design parameter, obtaining a $K^*$ that is logarithmic in the sample size \citep[Thm.~5.1, pp.~61475--61476]{zheng2024}, so it is not true that the thinning literature leaves the design unaddressed. The honest distinctions are three: their class is geometric, driven by a mixing time, where $\Bclass{r}$ is polynomial and the exponent is what the optimal design depends on; their coverage gap is marginal, where Definition~\ref{def:sharp-error} is calibration-conditional; and they prove no lower bound. The same three distinctions apply to the buffered three-way split of \citet{allohibi2026}, who work under geometric $\beta$-mixing and who themselves note that marginal validity follows from \citet{oliveira2024} and \citet{barber2026} without any buffer.

\paragraph{What this chapter claims.} Exactly five things. First, Theorem~\ref{thm:sharp}: the coefficient $1$ on the Kolmogorov drift term is sharp, by a construction that attains it, with the two correction terms identified and one of them shown exponentially small for split conformal. Second, Theorem~\ref{thm:block-b}: the exact minimiser of the expected calibration-conditional shortfall of the blocked route over $\Bclass{r}$, with its constant, and the statement that the balance point of the two terms is not the minimiser. Third, Theorem~\ref{thm:block-rate}: the blocked rate is not minimax in the calibration-conditional object at fixed confidence, with both halves proved, and the accounting of what blocking does buy. Fourth, Theorem~\ref{thm:block-lower}: the statement, construction and route of a mixing-driven lower bound at high confidence, together with the consistency check of Remark~\ref{rem:sharp-bp-check} and an explicit list of what is missing. Fifth, Proposition~\ref{prop:block-w1}: the specialisation of a published dependent Wasserstein bound to $\Bclass{r}$ with an explicit constant, claimed as a specialisation and nothing else. Theorem~\ref{thm:block-drift} is stated as a target with no part of it proved.

\section{Numerical illustration}\label{sec:sharp-example}

Two illustrations are given. The first sweeps the block length on the synthetic hedging-error process of Chapter~\ref{ch:certification}, to show what the design variable does to a certificate in practice. The second compares the blocked and unblocked bounds of Section~\ref{sec:sharp-notminimax} across record lengths, and has already been given as the table in Remark~\ref{rem:sharp-constants}.

\subsection{The block-length sweep}

The process is the one of Section~\ref{sec:cert-example} of Chapter~\ref{ch:certification}: a stationary Gaussian autoregression of order one for the per-period hedging shortfall, with $\varphi=\tfrac12$ and stationary variance one, a deployment mean shift of $\mu=0.05$ per period, $N=500$, $\alpha=0.1$, $\delta=\delta'=0.05$. Chapter~\ref{ch:certification} works the certificate at $b=10$. Table~\ref{tab:example-sweep} works it at three block lengths.

\begin{table}[htbp]
\centering\small
\caption{The certificate for the synthetic AR(1) example of Chapter~\ref{ch:certification}, $N=500$, $\alpha=0.1$, $\varphi=\tfrac12$, $\mu=0.05$, $\delta=\delta'=0.05$, at three block lengths. The coupling bound is the first inequality of Proposition~\ref{prop:prelim-ar1} of Chapter~\ref{ch:prelim}; $\dK$ is recomputed with $v_b$ at each $b$. At $b=5$ the coupling term dominates and the certificate is nearly vacuous; at $b=20$ it is negligible but $n$ has fallen to $12$ and the conformal term has grown.}\label{tab:example-sweep}
\begin{tabular}{@{}lrrr@{}}
\toprule
 & $b=5$ & $b=10$ & $b=20$ \\
\midrule
$n=\lfloor N/(2b)\rfloor$ & 50 & 25 & 12 \\
$k=\lceil0.9(n+1)\rceil$ & 46 & 24 & 12 \\
$k/(n+1)$ & 0.9020 & 0.9231 & 0.9231 \\
$\beta(b)$ bound & $1.563\times10^{-2}$ & $4.883\times10^{-4}$ & $4.768\times10^{-7}$ \\
$(n+1)\beta(b)$ & 0.7971 & 0.0127 & $6.2\times10^{-6}$ \\
$\dK(P,Q)$ & 0.0299 & 0.0391 & 0.0533 \\
Marginal level & 0.0730 & 0.8482 & 0.8467 \\
$b_{n,k}(0.05)$ & 0.8262 & 0.8239 & 0.7791 \\
$\beta(b)/\delta'$ & 0.3126 & 0.0098 & $9.5\times10^{-6}$ \\
Conditional level & 0.4837 & 0.7750 & 0.7258 \\
Calibration-side confidence $1-\delta-n\beta(b)$ & 0.1686 & 0.9378 & 0.9500 \\
\bottomrule
\end{tabular}
\end{table}

\datanote{The $b=5$ and $b=20$ columns are computed by the same formulas as the $b=10$ column of Chapter~\ref{ch:certification}: $v_5=5+2\sum_{h=1}^4(5-h)2^{-h}=5+2(2+0.75+0.25+0.0625)=11.125$, $\dK=2\Phi(0.25/(2\sqrt{11.125}))-1=2\Phi(0.03748)-1=0.0299$; $v_{20}=20+2\sum_{h=1}^{19}(20-h)2^{-h}=56.0000$ to four decimals, $\dK=2\Phi(1/(2\sqrt{56}))-1=2\Phi(0.06682)-1=0.0533$; the Beta quantiles are the roots of $\Prob(\mathrm{Bin}(50,u)\ge46)=0.05$, namely $u=0.8262$, and of $\Prob(\mathrm{Bin}(12,u)\ge12)=u^{12}=0.05$, namely $u=0.05^{1/12}=0.7791$. The last row is the calibration-side confidence $1-\delta-n\beta(b)$; the full conditional confidence of Theorem~\ref{thm:cov}(ii) subtracts a further $\delta'=0.05$, giving $0.1186$, $0.8878$ and $0.9000$, of which the middle value is the $0.8878$ reported in Chapter~\ref{ch:certification}. At $b=5$ the conditional level $0.4837$ is stated with calibration-side confidence only $0.1686$, because $n\beta(b)=0.78$; the column is reported to show the failure, not as a usable certificate. The AR(1) is normalised to stationary variance $v=1$, so the innovation standard deviation is $\sqrt{1-\varphi^2}$. Every cell is regenerated by the certificate package of Appendix~\ref{app:data}.}

\emph{Reading the table.} With $\varphi=\tfrac12$ the process forgets quickly and $b=10$ is already enough for the coupling term to be small; the certified level is then dominated by the drift and by the conformal term. For a slowly mixing process the picture changes. With $\varphi=0.8$, Proposition~\ref{prop:prelim-ar1} of Chapter~\ref{ch:prelim} gives $\beta(10)\le0.0538$ and $(n+1)\beta(10)\le1.40$: the certificate at $b=10$ is vacuous, and one must take $b=20$, where $n=12$, $\beta(20)\le0.00576$ and $(n+1)\beta(20)\le0.0749$, or else a longer record. The block length is not a nuisance parameter; it is the design variable of the method, and its choice requires a mixing class to be assumed.

\emph{What the sweep does not show.} A geometrically $\beta$-mixing autoregression lies in every class $\Bclass{r}$, because $|\varphi|^b$ is eventually smaller than $Cb^{-r}$ for every $r$. The sweep therefore shows only that some block length kills the coupling term and that too large a block length starves the block count. It does not locate $b^*$, whose value by Theorem~\ref{thm:block-b} depends on $r$ and would take the values in the table of Example~\ref{ex:sharp-bstar} on the same record. Locating $b^*$ numerically requires a process with genuinely polynomial mixing, and the natural candidate is the regime-renewal construction of Definition~\ref{def:sharp-renewal}, on which the exercise has not been carried out.

\figplan{Two panels for the synthetic AR(1) example, $\varphi\in\{0.5,0.8\}$. Horizontal axis: block length $b$ from $2$ to $50$ at $N=500$. Series: the coupling term $(n+1)\beta(b)$ from Proposition~\ref{prop:prelim-ar1} of Chapter~\ref{ch:prelim}; the conformal term $1-k/(n+1)$; the drift term $\dK(P,Q)$ at $\mu=0.05$; their sum; and the realised breach frequency of the certificate over $10^4$ Monte Carlo replications of the whole procedure, with the nominal $\alpha=0.1$ as a horizontal line.}{Deficit terms of the certificate as a function of block length for the synthetic hedging-error process. The coupling term falls geometrically in $b$, the conformal term rises as $n=\lfloor N/(2b)\rfloor$ falls, and the drift term rises with $b$ because the mean shift accumulates over the block. The Monte Carlo breach frequency tests the claim that the realised breach rate never exceeds $\alpha$ plus the sum of the three terms; the gap between the two curves measures the slack in the bound, most of it in the coupling term. Because the process is geometrically mixing and therefore lies in every class $\Bclass{r}$, the minimum of the summed curve is a property of this process and not an instance of the $b^*$ of Theorem~\ref{thm:block-b}.\label{fig:example-sweep}}

\datanote{The Monte Carlo series is produced by simulating $N_0+b+N+b$ periods of the AR(1), computing $\hat q$ from the blocked record, drawing the deployment block with the mean shift, and recording $\ind{S_0>\hat q}$; $10^4$ replications give a standard error of $0.003$ on a breach frequency near $0.1$. The realised breach frequencies at $b=5,10,20$ are $0.109$, $0.093$, $0.098$ for $\varphi=\tfrac12$ and $0.112$, $0.089$, $0.091$ for $\varphi=0.8$; each lies below the certified breach bound $\alpha+(n+1)\beta(b)+(1-k/(n+1))+\dK$, whose excess over $\alpha$ is $0.925$, $0.129$ and $0.130$ at $\varphi=\tfrac12$ and $8.72$, $1.50$ and $0.186$ at $\varphi=0.8$. Almost all the slack is in the coupling term. Code and seed are specified in Appendix~\ref{app:data}. No market data enter.}

\subsection{The design rule in use}

Collecting Theorem~\ref{thm:block-b}, Remark~\ref{rem:sharp-exact-b} and Remark~\ref{rem:sharp-constants}, the rule a bank should follow is short.

\begin{enumerate}[nosep,label=\arabic*.]
\item Assume a class: state $(C,r)$ with $r>1$, and record that Proposition~\ref{prop:noest} of Chapter~\ref{ch:certification} makes this an assumption about the market rather than an estimate from it.
\item Compute $b^*=(2(r+1))^{2/(2r+3)}N^{3/(2r+3)}$ from \eqref{eq:sharp-bstar}. Do not use the balance point $N^{3/(2r+3)}$, which costs between $27$ and $53$ per cent more over $r\in[1.5,5]$.
\item Evaluate the exact bound $g(b)=1/(2\sqrt{n+2})+(n+1)Cb^{-r}$ at the integers in a neighbourhood of $b^*$ and take the minimiser; the floor in $n$ moves it by a period or two, as Remark~\ref{rem:sharp-exact-b} shows.
\item Compare $\min_bg(b)$ with the unblocked Chebyshev bound $\sqrt{(\tfrac12+4C\zeta(r))/(N\delta)}$ of \eqref{eq:sharp-cheb} and use whichever is smaller, reporting both. Below about $10^5$ periods at $r=2$ the blocked bound wins; above it the unblocked sample quantile does.
\item If the blocked route is chosen, report the two-sided error of Proposition~\ref{prop:twosided}(c) of Chapter~\ref{ch:certification} at the chosen $b$, not the one-sided deficit, and state $(C,r)$, the drift assumption and $\delta,\delta'$ alongside it.
\end{enumerate}

\datanote{The range ``between $27$ and $53$ per cent'' in step~2 is the last row of the table in Example~\ref{ex:sharp-bstar}: $f(b_{\mathrm{bal}})/f(b^*)$ takes the values $1.275$, $1.327$, $1.411$ and $1.525$ at $r=1.5,2,3,5$. The crossover in step~4 is the table in Remark~\ref{rem:sharp-constants} and holds at $C=1$, $r=2$, $\delta=0.05$; it moves with all three.}

\section{Discussion}\label{sec:sharp-discussion}

\paragraph{What was shown.} Three statements are proved in full. The Kolmogorov drift charge of the certificate is exact: Theorem~\ref{thm:sharp} exhibits, for every calibration law with a continuous distribution function and every threshold procedure, a deployment law at Kolmogorov distance $\varepsilon$ whose coverage is lower by $\varepsilon$ up to two corrections, one removable by taking the far point far and one exponentially small in the number of blocks for split conformal. The block length of the blocked route is optimised exactly over a polynomial mixing class: Theorem~\ref{thm:block-b} minimises the expected calibration-conditional shortfall implied by the Beta law and the coupling term, gives the minimiser and its constant, and shows that the point where the two terms balance is not the minimiser but costs a third more at the parameters of the worked example. And the resulting rate is not minimax: Theorem~\ref{thm:block-rate} proves that the unblocked sample quantile of the whole record attains $(N\delta)^{-1/2}$ over every class $\Bclass{r}$ with $r>1$, by a covariance bound on half-line indicators, the variance of the empirical distribution function and Chebyshev's inequality applied at two quantiles, and that no procedure beats $N^{-1/2}$ at fixed confidence, by a two-point argument on independent scores which lie in every class.

\paragraph{What blocking buys, stated once.} It does not buy rate. It buys three things, and each is used in Chapter~\ref{ch:certification}: the exact $\Beta(k,n+1-k)$ law of the calibration-conditional coverage, hence an exact quantile in place of a concentration inequality; a drift term charged on a half-line that the coupling has frozen, hence the constant $1$ that Theorem~\ref{thm:sharp} shows to be sharp; and a certificate whose only dependence on the assumed mixing class is the single term $(n+1)\beta(b)$, so that a reader who disputes the class can recompute the level in one multiplication. Remark~\ref{rem:sharp-constants} adds a fourth, which is numerical rather than structural: at every record length a bank possesses, the blocked bound is smaller than the proved unblocked bound, because the factor $\sqrt{b^*}$ grows slowly while the constant $\sqrt{(\tfrac12+4C\zeta(r))/\delta}$ is paid at once. The crossover is near $10^5$ periods at $r=2$.

\paragraph{What is not shown.} Four things, and they are stated in the boxes rather than in prose. The mixing-driven lower bound at high confidence, Theorem~\ref{thm:block-lower}, is a statement with a construction and a route: its part~(a) needs the concentration of an excised record, its part~(b) needs the two-point reduction over all procedures, its Lemma~\ref{lem:sharp-renewal-beta} needs the constant relating the $\beta$-coefficient to the no-renewal probability, its $\delta$-dependence $g_r$ is open between $\sqrt{\log(1/\delta)}$ and $\delta^{-1/r}$, and its Fuk--Nagaev companion rests on an unverified citation. The partially observed drift problem, Theorem~\ref{thm:block-drift}, has no part proved. The Wasserstein specialisation, Proposition~\ref{prop:block-w1}, carries an unextracted absolute constant and is in any case a specialisation of somebody else's theorem. And the constants of Theorem~\ref{thm:block-rate} differ by a factor of $316$ at the parameters of the example, so only the exponent is matched.

\paragraph{Limitations of the whole approach.} Three are worth naming beyond the individual gaps. First, everything here is conditional on an assumed mixing class, and Proposition~\ref{prop:noest} of Chapter~\ref{ch:certification} shows that the class cannot be established distribution-free from any finite record; the optimal block length is optimal for the class one assumes, and a process outside it defeats the design as thoroughly as it defeats the certificate. Second, the error functional of Definition~\ref{def:sharp-error} strips out the drift and the deployment coupling, which is legitimate because both are additive and because Theorem~\ref{thm:sharp} shows the drift charge cannot be improved by design, but it means the optimisation is of a part and not of the whole; the whole is the two-sided bound of Proposition~\ref{prop:twosided}(c) of Chapter~\ref{ch:certification}, and at $N=500$ the drift and Markov terms together contribute $0.049$ of the $0.1345$, so a third of the error is outside the optimisation. Third, no illustration of $b^*$ on a genuinely polynomially mixing process has been run: the sweep of Section~\ref{sec:sharp-example} is on a geometrically mixing autoregression, which lies in every class, and the honest statement is that the design rule has been derived and not yet demonstrated.

\paragraph{What carries doctoral weight.} This should be said plainly rather than left to inference. Theorem~\ref{thm:block-b} is elementary. It is one differentiation of a two-term function, and its value lies in the exactness of the constant and in the correction of the natural but wrong instinct to balance the two terms, not in any difficulty. Theorem~\ref{thm:sharp} is a short construction; it is the kind of result a referee expects to see beside a bound of the form of Theorem~\ref{thm:cov}, and its absence would be a defect, but its presence is not a contribution of scale. Proposition~\ref{prop:block-w1} is a specialisation and is claimed as one. Theorem~\ref{thm:block-drift} is a plan.

Two things do carry weight, and one of them is negative. The negative one is Theorem~\ref{thm:block-rate}. The literature on conformal prediction under dependence, this thesis included at an earlier stage of its own argument, treats blocking as the route to a guarantee and does not ask what the guarantee costs relative to doing nothing. The answer is a factor $\sqrt{b}$ in effective sample size, and the honest consequence is that the blocked rate $N^{-r/(2r+3)}$ is not minimax in any coverage formulation: marginally it is dominated by \citet[Cor.~2]{barber2026} without blocking, and conditionally it is dominated by the ordinary sample quantile for every $r>1$. Establishing that, with both halves proved and with the numerical crossover computed so that a practitioner knows which route to use at the record length in hand, is a real result about a design that is in use. It is also the result that determines what the positive claim of the chapter has to be.

The positive one is Theorem~\ref{thm:block-lower}, and it is not proved. The chapter's doctoral claim therefore rests on work still to be done, and the honest position is to say so and to say what the work is. It has three parts: the regime-renewal process of Definition~\ref{def:sharp-renewal} must be shown to lie in $\Bclass{r}$ with an explicit constant; the displacement argument of Remark~\ref{rem:sharp-run} must be made a proof for the sample quantile; and the two-point reduction must be carried out over all threshold procedures. The first is routine renewal theory. The second is a concentration argument on a conditioned, non-stationary remainder. The third is the hard one, and it is the one that answers the question \citet[Sec.~4]{barber2026} leave open, in its calibration-conditional form. Remark~\ref{rem:sharp-bp-check} gives the one piece of external evidence available: the claimed high-confidence rate integrates over the confidence level to $N^{-r/(r+1)}$, which is the rate they prove for split conformal marginally, so the construction is at least calibrated against a known answer. If the reduction holds only over the regime-renewal subclass, or loses a logarithmic factor, the statement is weaker and still says what it needs to say, namely that the mixing exponent, and not the Gaussian scale, sets the achievable calibration-conditional coverage error at the confidence levels a certificate is written at. If it fails altogether, the chapter retains Theorems~\ref{thm:sharp}, \ref{thm:block-b} and \ref{thm:block-rate}, and the thesis retains a certificate that is sharp in its drift term, optimally designed within its route, and honest about the route's cost. That is the floor, and it is stated as the floor.

\paragraph{Connection to the empirical chapters.} Chapter~\ref{ch:usd} runs the certificate on the \gls{libor}-to-\gls{sofr} record, where the deployment law became observable and the assumed class can be tested indirectly by comparing realised breach frequencies with certified levels. The design rule of Section~\ref{sec:sharp-example} fixes the block length there, and Theorem~\ref{thm:block-rate} fixes what else must be reported alongside it: the unblocked bound, so that the reader can see which of the two routes the record length favours. Chapter~\ref{ch:zar} applies the same rule to the converted \gls{jibar} book with the short \gls{zaronia} record available, where $N$ is small, the blocked bound is the smaller of the two by a wide margin, and the assumed class is doing most of the work. Both chapters report the two-sided error of Proposition~\ref{prop:twosided}(c) of Chapter~\ref{ch:certification} at the chosen block length, which is the quantity this chapter has been about.

\section{Chapter notes}\label{sec:sharp-notes}

The device of moving the extremal mass of a law to a far point, to make a distributional-distance bound tight, is standard in robust statistics and is used here in the form dictated by the Kolmogorov ball; the specific observation that the mass must come from the \emph{bottom} of the calibration law, and that taking it proportionally from the whole law as a mixture attains only the constant $1-\alpha$, is Remark~\ref{rem:sharp-mixture} and is the point at which a plausible first attempt fails.

The covariance bound of Lemma~\ref{lem:sharp-cov} is the definition of the strong mixing coefficient and needs no citation beyond Definition~\ref{def:prelim-ab-sigma} of Chapter~\ref{ch:prelim}; the resulting variance bound of Lemma~\ref{lem:sharp-var} is the textbook estimate for the empirical distribution function of a mixing sequence and the general reference is \citet{doukhan1994}. The two-point method of Theorem~\ref{thm:block-rate}(b) is Le Cam's, in the form that combines a separation condition with Pinsker's inequality \citep[\S2.3, \S2.4.1]{tsybakov2009}; the version stated there with the constant $\tfrac14e^{-\mathrm{KL}}$ is not used, because the total-variation form gives a cleaner constant and needs no more than the definition.

The regime-renewal construction of Definition~\ref{def:sharp-renewal} is a semi-Markov process with heavy-tailed sojourns, and the observation that such a process has polynomial rather than geometric mixing is classical. Its use as a lower-bound construction for a conformal prediction problem is new so far as the searches recorded in this thesis go; \citet[Ex.~2, p.~12]{oliveira2024} give a $\beta$-mixing example built from a random walk on a cycle graph, whose coefficients decay geometrically and which therefore lies in every class $\Bclass{r}$, so it cannot serve. The hidden-copy process of Proposition~\ref{prop:noest} of Chapter~\ref{ch:certification} shares the mechanism of a rare event that is invisible in the record, and the coupling step of its proof is the model for the route of Lemma~\ref{lem:sharp-renewal-beta}.

Two further remarks on the literature. The bracket \eqref{eq:sharp-minimax} has the same shape as the four-term decomposition of \citet[Thm.~2]{halkiewicz2026}, with his noise term $\sqrt{1/m}$ playing the role of the Gaussian term and his drift term $(m/T)^\beta$ playing the role of the mixing term; the shapes agree because both are the result of balancing a concentration term against a bias term, and the mathematics differs entirely in what produces the bias. And the conditional object of Definition~\ref{def:sharp-error} is the one that \citet[Sec.~2.2]{zwart2025} report as already in operational use under exchangeability through the Beta quantile; what is new here is the optimisation of a design parameter against it under dependence, not the object.
